% UNINTEGRATED WHOLE-B18 WORKING APPENDIX, resumed 25 October 2026.
% Same existing general compact paper family; not a standalone paper.
% Complete source transcription at the stated shared-appendix interfaces.
% DRAFT: independent source review, integration and compilation are pending.
% No compilation, worker import, test or scientific execution has occurred.

\section{Terminal walls in a whole \(B_{18}\) diagram}
\label{app:wholeB18}
Fix the original strict interior unit point \(z\) and an original
farthest vertex \(p\), and assume that the entire incident diagram
is \(B_{18}\). Normalize the actual old normals by
\[
 n_i=(e_i-e_{i+1})/\sqrt2\quad(1\leq i\leq17),\qquad n_{18}=e_{18}.
\]
Write
\[
 p^2=-1,\qquad z=E(p-x),\qquad E>0,\quad x\perp p,\quad
 x_1>\cdots>x_{18}>0,\quad \norm x^2<1.
\]
For \(P_i=e_1+\cdots+e_i\), put
\[
 S_i=\ip x{P_i},\qquad f_i=\sqrt i-.58/\sqrt2,\qquad
 d_K=\sqrt K-\sqrt{K-1}.
\]
The actual duals are \(\omega_i=\sqrt2P_i\) for \(i\leq17\)
and \(\omega_{18}=P_{18}\). In particular
\(\ip{\omega_i}{\omega_j}=2\min(i,j)\) for \(i,j\leq17\).
An index is strong when \(S_i>f_i\); equality belongs to the weak case.
There are no old spectator directions. All subsequent references
in the original frame concern this one vector \(x\).

We shall prove that these hypotheses are impossible. The proof first
constructs unconditional terminal sum bounds, then classifies the actual
seventeenth opposite. Its far-helper, near-six and short-contact branches
are treated by shared actual-point and neighboring-comparison arguments.
No bound on an arithmetic field occurs.

We use the actual positive projection and future-point pressure lemmas
already proved in the main text. Their hypotheses will be checked for
each point below. A positive correction on actual old normals leaves
every unknown wall inequality nonpositive by acuteness. A spacelike
corrected normal is used only inside its positive finite projection;
it is never normalized as a chamber point.

\subsection{The early window and actual head pressures}
\label{app:B18bulk}
Deleting \(i\), \(5\leq i\leq14\), retains
\(A_{i-1}+B_{18-i}\), with both ranks at least four.
The finite one-node list consists of isolation; one left end bond
of label three or four; a left \(D\) or permitted \(E\) position;
one right long-end simple bond; or an end-to-end bridge.
Two contacts on one retained component create a cycle. A joined
rank-eighteen diagram is classical; exceptional types add no case.
The right \(B_3\) at \(i=15\) additionally allows its short-end
simple bond, and the right \(B_2\) at \(i=16\) allows either end.
These last roles will be checked separately.

For an actual edge with \(C=\cosh L=qR\), \(q=\sqrt{2i}\),
clearing the positive denominators of its directional pressure
shows that
\[
 R>T(q):=\frac{50}{29}
    -\frac{q-29/50}{q(2q-29/50)}
 \quad\Longrightarrow\quad S_i>f_i.
\]
The function increases for \(q\geq\sqrt{10}\), and substitution gives
\[
 T(q)<83/50\ (q<7),\qquad
 T(q)<801/500\ (q<15/4),\qquad
 T(q)<159/100\ (q<16/5).
\]
The finite-white Schur calculations have the following complete
non-helper bounds; \(\phi=(1+\sqrt5)/2\).
\[
\begin{array}{c|l}
\text{role}&\text{lower bound for }R\\ \hline
\text{left near three, own at least five}
 &[i(\phi+1)-1]/\sqrt{2i(i+1)}\\
D\text{ position}&(i\phi/2+1)/\sqrt{2i}\\
E\text{ position}&2\phi/2+3/8\\
\text{right long three, positive remaining own}
 &1+1/\sqrt2\\
\text{right }B_3\text{ short three, positive own}
 &1+1/\sqrt2\\
\text{end-to-end bridge}&(i+1)/\sqrt{2i}.
\end{array}
\]
The near-three bound increases with \(i\). At \(i=6\) it is
\((6\phi+5)/\sqrt{84}>801/500\), which covers \(6,7\);
at \(i=8\) it exceeds \(83/50\), covering larger \(i\).
At \(i=5\), every own at least seven gives \(R>5/3\).
For the \(D\) row, \(i=5\) gives \(R>159/100\), and
\(i\geq6\) gives \(R\geq2927/1750>83/50\).
An \(E\) position has \(i\leq8\) and retained residual at most
\(1/4\). The bridge residual is \(1/(2i)\); its numerator is
at least \((i+1)/(2i)\), proving the displayed lower bound,
including own zero. Left label-four roles have the affine
small-own alternative or the larger direct edge bounds:
near own three has \(C=(i-1)+i/\sqrt2\).
The positive-null and impossible-time boundaries in the common
prefix inventory apply to the actual retained blocks at these
ranks. For the right long-three own-zero boundary one can also
see the kernel directly: use coefficient one on the cap and
old \(i\), two on old \(i+1,\ldots,17\), and \(\sqrt2\) on
old short \(18\). Every selected product vanishes.

An isolated finite cap has the actual finite/core pressures used
in the earlier component lemmas. At own five its loss is less than
\(25/76+5/24<.58\); own seven has loss
\(1000/4779+9/25<.57\); the remaining finite interval has its
largest loss at \(i=5\), own eight:
\(\sqrt{10}+\sqrt2-4<.578\).
For an isolated dotted cap put \(u=g^{-2}\). Its positive
reciprocal core gives
\[
 \ip x{\omega_i}>\sqrt{2i-u}-\sqrt{u(1-u)}.
\]
For \(i\geq6\) the loss is at most
\[
 a_i u+\sqrt{u(1-u)},\qquad
 a_i=\frac1{\sqrt{2i}+\sqrt{2i-1}}<37/250,
\]
whose maximum is less than \(.5795\), since
\(\sqrt{1+(37/250)^2}<1011/1000\).
Thus only an isolated dotted cap at \(5\) may remain weak.

There is one early near-three exception, own five at \(5\).
For \(g=\phi/2\), the actual cap has displacement
\(v=g\omega_5+\omega_4/2\). On the right actual \(H_4\), correct by
\[
 w=g(3n_5/2+n_6+n_7/2),\qquad
 \delta=(5-3\phi)/8>0.
\]
The corrected normal is spacelike and its finite projection is
actual; the extra old \(4\) product is negative, and all unknown
wall signs follow by positivity. Its pure spatial displacement is
\[
 Y=(1/2+3g/4)\omega_4+(g/4)\omega_8,\qquad
 A^2=(7+13\phi)/2,\quad Y^2=(33+49\phi)/8.
\]
The finite-point pressure gives
\[
 \ip xY>18/5,\qquad Y^2<281/20,\qquad
 \sum Y_j<43/5.
\]
Independently, the earlier rank-valid \(.61\) prefix lemma supplies
a strong \(K=11\) or \(12\). Its old floors exceed \(72/25\)
or \(303/100\). The vector \(Y+(3/4)P_K\) therefore has
squared demand exceeding squared norm by \(.0401\) or
\(.78625625\), respectively, a contradiction.
This use of the earlier lemma concerns \(B_{18}\geq B_{16}\)
and precedes any new \(.58\) or strong-five conclusion.

The two extra \(B_2\) roles at \(16\) deserve separate treatment.
For short-three/own-five put \(g=\phi/2\) and
\(\alpha=\sqrt2g+1/2\). The left actual \(H_4\), on the cap and
old \(16,15,14\), has residual \((5-3\phi)/8\), time square
\(16\alpha^2-1/2\), and normalized spatial displacement constant
through \(13\), of norm less than \(4\). Since \(\alpha>1.643\),
its pressure exceeds \(3.976-.083=3.893>3.89\).
The earlier, independently available strong \(K=11,12\) references
have energies \(>.8472225,>.8689525\), sums \(<2.886,<3.034\),
and terminal capacities \(>1/7\). Their roots and actual points
lie through \(14\). Paying \(1/7\) at that terminal anchor gives
\[
 .8472225+2(3.89-2.886)/7-5/49>1.0320,
 \qquad
 .8689525+2(3.89-3.034)/7-4/49>1.0318.
\]
For short-three/own-three the same cap has displacement
\(v=(\omega_{16}+\omega_{18})/2\). Deleting actual short \(18\)
instead gives a positive full \(D_{18}\) projection. Its sum and
time square are \(9+8\sqrt2\) and \(23/2+8\sqrt2\); the squared
sum minus eighteen times the time square is \(2\). Thus its actual
short-edge cosine is \(8+9\sqrt2/2>14\), whence \(S_{18}>3.949\).
The unspent early \(K=14,15\) references, of effective energies
\(>.9231,>.94006\), have terminal capacities \(>1/9\).
Paying \(1/9\) toward \(P_{18}\) gives respectively
\[
 .9231+2(3.949-3.332)/9-4/81>1.0108,
 \quad .94006+2(3.949-3.463)/9-3/81>1.0110.
\]
The \(H_4\) construction is a positive finite projection; the second
point is the positive projection onto the displayed full finite Gram.
All original unselected wall products are therefore nonpositive.
Neither new point is asserted farthest. These arguments use the
early references before the new \(B_2\) conclusion, so they introduce
no circularity.

Consequently every weak ordinary site is a pure far helper
\[
 c_i=A_i p-g\sqrt2P_i-e_1/\sqrt2,\qquad
 A_i^2=2ig^2+2g-\tfrac12,\qquad
 \phi/2\leq g<6/5,
\]
apart from the isolated dotted \(5\).
The upper cutoff follows from
\(R^2\geq4g-2\geq14/5>(83/50)^2\) at \(g\geq6/5\).

For adjacent pure helpers the exact time/dual calculation gives
\[
 \ip{c_i}{c_{i+1}}
 \geq \frac12-g_i g_{i+1}
           +\frac14(g_i/g_{i+1}+g_{i+1}/g_i)
 \geq1-g_i g_{i+1}>-11/25.
\]
Their actual native gap forces orthogonality. The positive two-cap
replacement, with every other old wall retained, has positive
inverse Schur matrix; its weighted loss is less than
\(7/16+1/(4\sqrt{10})<.58\).
Thus both sites cannot be weak. This uses their actual new
parameter range, rather than inclusion in an older weak class.
In particular some \(K=14\) or \(15\) is strong.

The same actual helper feet give \(x_1>.5455\) at a weak
helper \(5\). For an isolated dotted weak \(5\), its reciprocal
core gives
\[
 U'=P_5+(u/2)(e_6-e_5),\qquad
 \ip x{U'}>T'>1.8185,\qquad .4<u<1.
\]
Indeed the loss is bounded by
\(a u+\sqrt{u(1-u)}\), \(a=\sqrt{10}-3<13/80\);
for \(u\leq.4\) it is less than \(.565<.58\), and globally
it is less than \(.588\).
The pressure bound follows from
\(\sqrt{10}>3.162,\sqrt2<99/70\).
These are actual point pressures, not assigned floors at weak nodes.

\subsection{Raw references, head isolation and the uniform paid prefix}
\label{app:B18references}
A reference \(w\) has support through a strong \(K\), sum \(f_K\),
actual pressure \(\ip xw>A_w\), and effective energy
\(L=2A_w-w^2\). Point insertion need not leave \(L=w^2\).
For a new actual pressure \(\ip xV\geq G\), weight \(\mu\geq0\),
and \(Y=V-P_K\), the completed-square rule is
\[
 \norm x^2>L+2\mu(G-f_K-\ip wY)-\mu^2Y^2,
 \qquad \mu\leq\beta_K .
\]
The inequality is strict when the old strong reference retains
a positive coefficient. Every use below checks its individual prices.

Let \(d_i=\sqrt{i+1}-\sqrt i\).
The first-five energy exceeds \(.6668157\), and rational lower
numerators for \(10^7d_i^2\), \(i=5,\ldots,16\), are
\[
 455488,385186,333704,294372,263340,238230,
 217494,200080,185248,172465,161332,151549.
\]
For a direct verification, put \(a=4i+2\) and use
\(d_i^2(a-d_i^2)=1\), whence
\(d_i^2>(a^2-1)/(a(a^2-2))\).
Thus the unomitted energies satisfy
\[
 E_{14}>.9241299,\quad E_{15}>.9413764,\quad
 E_{16}>.9575096,\quad E_{17}>.9726645.
\]
A missing knot \(j\geq6\) loses less than
\(1/[32(j-1)^3]\); nonadjacency makes these losses disjoint.
Their total is less than
\((1/32)\int_4^\infty t^{-3}\,dt=1/1024\).
Each terminal capacity is at least \(d_K\).

Before any head payment, the three initial-five cases give the
following full contracts:
\[
\begin{array}{c|cc|c}
\text{case at five}&L_{14}>&L_{15}>&w_1<\\ \hline
\text{strong}&.9231&.9403998&.3652\\
\text{pure helper}&.95255&.96895&.5455\\
\text{isolated dotted weak}&.9291&.9455&.377 .
\end{array}
\]
In the strong case interpolate through all actual strong floors,
skipping every weak knot. The first mean is \(f_5/5<.3652\).
For a helper, interpolate through \((1,.5455),(6,f_6),(K,f_K)\).
Its first-six energy exceeds
\(.5455^2+(f_6-.5455)^2/5>.74385\);
the tails to \(14,15\) exceed \(.2087,.2251\).
All slopes are positive decreasing.

For isolated dotted weak \(5\), choose strong \(J=6\) or \(7\)
and a reference constant first \(J\), then straight to \(K\).
Insert \(\mu[U'-(5/J)P_J]\), with \(\mu=.21,.2\).
The residual squares are bounded by
\[
 5-u+u^2/2-25/J<77/150,\ 194/175
\]
and the deficits exceed \(.119,.221\).
Its individual \(P_J\) prices are \(.175,1/7\),
below the capacities \(>.1783,>.1628\) already at \(K=14\).
The gains exceed \(.027342,.044057\).
The bare first-\(J\) energies exceed \(.69312,.7139\);
the tails to \(14\) exceed \(.2087,.1715\) and to \(15\)
exceed \(.2251,.18825\).
The full first coordinates are less than
\[
 .3399+.21(1-5/6)=.3749,\qquad
 .31938+.2(1-5/7)<.377.
\]
This uses the straight reference for these prices. A curved
reference with a smaller \(P_6\) capacity is not substituted.

The actual long-head opposite is isolated. Deleting old \(1\)
retains \(B_{17}\); its finite enlargement is isolation or a
simple bond to its long endpoint \(2\). In the latter case
\[
 v=g\omega_1+\omega_2/2,\qquad A^2=2g^2+2g,\qquad C_1=2g+1.
\]
Own zero has zero time; own three is an actual affine triangle.
Every other own gives \(x_1>\sqrt{\sqrt2-1}>.643\).
If \(5\) is strong, replace the first-five mean \(<.3652\)
by this pressure, gaining
\((5/4)(.643-.3652)^2>.0964\) over \(L_{14}>.9231\).
If \(5\) is weak, choose strong \(J=6,7\) and interpolate
through \(1,J,K\); the first energies exceed \(.8034,.8361\)
and the tails exceed \(.2087,.1715\).
Each case exceeds one, including isolated dotted weak \(5\).
Therefore a head has time \(A_H^2=2g_H^2-1\).
Owns zero, three, four have impossible time; six and parallel
have the positive actual \(6/3\) and parallel nulls.
Only own five, finite own at least seven, and dotted own remain.

For later reference, every non-five head has a point pressure
\(\ip xZ\geq h\), with \(\sum Z_j=1\) and the following
full local data:
\[
\begin{array}{c|c|c|c}
\text{own}&Z&h&R=Z^2\\ \hline
\text{finite }\geq7&
 (1-u/2)(e_1+e_2)+(u-1)e_3&
 [\sqrt{2-u}-\sqrt{7u-4-3u^2}]/\sqrt2&
 3-4u+3u^2/2\\
\text{dotted}&(1-u/2)e_1+(u/2)e_2&
 [\sqrt{2-u}-\sqrt{u(1-u)}]/\sqrt2&
 1-u+u^2/2 .
\end{array}
\]
Here \(u=g_H^{-2}\); the first row has
\(1<u\leq\cos(\pi/7)^{-2}\), the second \(0<u<1\).
Their actual positive cores are
\(c_H+g_H^{-1}n_1+2(g_H^{-1}-g_H)n_2\)
and \(c_H+g_H^{-1}n_1\), respectively.
Their cap product is zero, the stated old-root products are
nonpositive and all unknown signs follow termwise.
They are future, giving these pure spatial pressures.
Both obey
\begin{equation}\label{app:eq:B18headF}
 .8h-.16R>.2893.
\end{equation}
For native seven, \(.9009<g_H<.901\) gives \(h>.4315,R<.349\),
and the left side exceeds \(.28936\).
For finite labels at least eight, \(1<u<1.172\).
Writing \(D=(u-1)(4-3u)\), the derivative is
\[
 \frac{.8}{\sqrt2}\left(-\frac1{2\sqrt{2-u}}
                 -\frac{7-6u}{2\sqrt D}\right)+.64-.48u.
\]
It is negative on \((1,7/6]\); on \([7/6,1.172)\),
\(\sqrt D>.284\), its positive term is \(<.032\)
and its other terms total less than \(-.20\).
At own eight \(h>.4388,R<.373\), giving \(>.29136\).
For dotted \(u\), the two chords
\(\sqrt{2-u}\geq1.4142-.379u\) on \([0,1/2]\) and
\(\sqrt{2-u}\geq1.4494-.4494u\) on \([1/2,1]\)
give, respectively,
\[
 .8h-.16R>
 .63987152-.0944u-.5657\sqrt{u(1-u)}>.30587152,
\]
\[
 .8h-.16R>
 .69978064-.2142u-.5657\sqrt{u(1-u)}>.29023064.
\]
Use
\(bu+c\sqrt{u(1-u)}\leq(b+\sqrt{b^2+c^2})/2\),
with radical bounds \(.5736,.6049\).
This proves the entire finite and unbounded dotted head seal.

\subsection{One uniform paid prefix, without a terminal premise}
\label{app:B18paidprefix}
If the head has own five, its actual unit normal is
\[
 c_H=A_Hp-(\phi/\sqrt2)e_1,\qquad A_H^2=(\phi-1)/2.
\]
A weak pure helper at \(i=5,6,7\) would have mutual magnitude
\[
 M_i(g)=A_H\sqrt{2ig^2+2g-\tfrac12}
             -\frac\phi2(2g+1).
\]
This increases in \(g,i\) on the allowed range and
\(M_7(6/5)<0\), since the positive squared comparison has
difference \((814\phi-1392)/100<0\).
An isolated dotted weak \(5\) has own \(h<2\), by its edge
pressure and \(f_5^2<7/2\). Its increasing mutual magnitude
\[
 A_H\sqrt{10h^2-1}-\phi h
\]
lies strictly between zero and \(1/4\) for \(1<h<2\):
the endpoint squared comparisons are
\((7\phi-11)/2>0\) and \(377/16-(29/2)\phi>0\).
Both cases contradict actual acuteness or the native gap.
Thus a native-five head makes \(5,6,7\) strong independently
of any terminal hypothesis.

The head's actual finite \(H_4\) projection gives
\[
 \ip x{P_4/4}\geq h_0,\qquad
 h_0^2=(\phi-1)/4,\qquad h_0>.39307.
\]
At strong \(K=14,15\) the curved reference, omitting every
nonadjacent internal weak knot, has energies
\[
 w_{14}^2>.924016575,\qquad w_{15}^2>.94125107.
\]
Indeed the complete raw energies above lose at most
\((.0112^2+.008^2+.0061^2)/2=.000113325\)
for \(K=14\), and \(<.00012533\) for \(K=15\);
these bounds include every permitted hole.
The \(P_5\) capacity is \(>.151\).
Pay weight \(.56\) on \(P_4/4-P_5/5\), costing \(.112\).
Its square is \(1/20\), it is orthogonal to the raw reference,
and the single-square gain exceeds
\[
 1.12(.39307-.3652)-.56^2/20=.0155344.
\]
Hence this initial/head case gives
\[
 L_{14}>.939550975,\qquad L_{15}>.95678547.
\]

All other initial/head cases give larger uniform energies.
For a non-five head and strong \(5\), pay \(.4(Z-P_5/5)\).
Its price is \(.08<.151\), and
\eqref{app:eq:B18headF} gives gain
\(.2893-.8(.3652)+.032>.02914\).
For a weak helper \(5\), the raw contracts above already
give \(L_{14}>.95255,L_{15}>.96895\), with no further head graft.

For isolated dotted weak \(5\), choose strong \(J=6,7\) and
a curved reference to \(K\), constant first \(J\).
Its first mean, individual capacity, and the square of
\(Y_5=U'-(5/J)P_J\) obey
\[
\begin{array}{c|ccc}
J&m<&\beta_J>&Y_5^2<\\ \hline
6&.3399&.1435&77/150\\
7&.31938&.1366&194/175 .
\end{array}
\]
The raw first-\(J\) energies exceed \(.69317,.71400\);
the complete tails to \(14\) exceed \(.2117654,.1732468\)
after omission losses \(<.000267,.000181\).
Thus the raw energies exceed \(.9046684,.8870658\).
At \(15\), the additional increment is \(>.0172465\).
Let \(Y_H=Z-P_J/J\). The full cross is
\(\ip{Y_H}{Y_5}=1-5/J\), while
\(\ip w{Y_H}=\ip w{Y_5}=0\).
The one vector \(w+.4Y_H+.09Y_5\) spends \(.85/J\),
below each displayed capacity. Its entire gain is
\[
 (.8h-.16R)-.8m+.16/J
 +.18(T'-5m)-.0081Y_5^2-.072(1-5/J).
\]
It exceeds \(.04930,.06699\), respectively. The resulting
\(K=14\) energies exceed \(.9539684,.9540558\), and the
\(K=15\) energies exceed \(.9712149,.9713023\).
The earlier \(.21/.2\) insertion is not also charged.

In every case, therefore, there is one paid reference \(V\),
supported through a strong \(K=14\) or \(15\), with sum \(f_K\),
unspent terminal capacity \(\beta_K\geq d_K\), and
\begin{equation}\label{app:eq:B18uniformprefix}
 2A_V-V^2>L_K,\qquad
 L_{14}=.939550975,\quad L_{15}=.95678547 .
\end{equation}
All corrections have zero sum and support through seven.
Their individual nonnegative coefficients were checked before
the terminal pressure is used.

We shall also need the following independent implication:
if \(17\) is strong, the head has own five.
For a non-five head and strong \(5\), the curved reference to
\(17\) has energy \(>.97133\), and its head gain \(>.02914\)
already exceeds one. If \(5\) is a helper, the full curved
reference gains \(>.74385-.7126=.03125\) in its first six
coordinates, yielding
\(E_{17}+.03125-1/1024>1.0029\).
If \(5\) is isolated dotted, the analogous curved references
to \(17\) have energies \(>.9530,.9354\); the joint
head/isolated-five gains just proved exceed \(.04930,.06699\).
Again the totals exceed one. Thus every non-five head is
incompatible with strong \(17\). This argument precedes any
classification of its terminal opposite.

\subsection{Unconditional terminal sums and isolated opposites}
\label{app:B18sumcutoffs}
For \(N=16,17,18\), let \(Y=P_N-P_K\).
It has support after \(K\), square \(N-K\), and zero product
with the entire corrected reference in
\eqref{app:eq:B18uniformprefix}.
Paying the nonnegative weight \(d_K\leq\beta_K\) shows that
\(S_N\geq G\) would imply
\[
 \norm x^2>L_K+2d_K(G-f_K)-(N-K)d_K^2.
\]
The rational root bounds
\[
\begin{gathered}
3.60555<\sqrt{13}<3.60556,\quad
3.74165<\sqrt{14}<3.74166,\\
3.87298<\sqrt{15}<3.87299
\end{gathered}
\]
give
\(.1361<d_{14}<.1362,\ .1313<d_{15}<.1314\).
Every bound follows by positive squares and
\(d_K=1/(\sqrt K+\sqrt{K-1})\).

For \(S_{17}\geq3.760\), the extra terms are bounded below by
\[
 2(.1361)(.428)-3(.1362)^2=.06085028,\qquad
 2(.1313)(.297)-2(.1314)^2=.04346028.
\]
Their totals exceed \(1.000401255,1.00024575\).
For \(S_{18}\geq3.830\), the increments are
\(.06135384,.04457632\), giving totals
\(>1.000904815,>1.00136179\).
Thus
\begin{equation}\label{app:eq:B18sumbounds}
 S_{17}<3.760,\qquad S_{18}<3.830,\qquad
 C_{17}<10.879,\quad C_{18}<9.81 .
\end{equation}
The edge cosines use the actual directions
\(P_N/\sqrt N\):
\(C_N<(N+G_N^2)/(N-G_N^2)\).
The two fractions are
\(31.1376/2.8624<10.879\) and
\(32.6689/3.3311<9.81\).
These are simultaneous consequences for the same \(x\),
not energies added from different references.

An entirely isolated \(17\)-opposite has
\[
 c=A p-\sqrt2gP_{17},\qquad A^2=34g^2-1>0.
\]
Own zero has negative time. Own three, four, six and parallel
have positive actual null combinations, respectively,
\[
 c+n_{16}+2n_{17}+\sqrt2n_{18},\qquad
 c+\sqrt2n_{17}+n_{18},\qquad
 g c+n_{17}+n_{16}/2,\qquad c+n_{17}.
\]
Selected products vanish, cap time is positive, and all
unselected wall products are nonpositive. Compactness forbids them.
For own five, correct on the actual \(H_4\) by
\(g(3n_{17}/2+n_{16}+n_{15}/2)\), with positive residual
\(\delta=(5-3\phi)/8\). The actual finite projection gives
pressure on \(V=P_{14}+\tfrac34(e_{15}+\cdots+e_{18})\)
strictly greater than \((A-\sqrt\delta)/(\sqrt2g)\).
Ordering implies \(\ip xV\leq S_{17}\).
Since \(A^2=(17\phi+15)/2\), the rational bounds
\(A>4.61,\sqrt\delta<.136,\sqrt2g<1.145\) give
\(S_{17}>(4.61-.136)/1.145>31/8\).
Every finite own at least seven has the positive core
\(c+g^{-1}n_{17}+2(1-g^2)g^{-1}n_{16}\).
Its norm is \(-(1-g^2)(4g^2-3)/g^2\), selected cap/own
products vanish, and the other selected product is negative.
All unknown wall signs follow from positive actual coefficients.
With \(u=g^{-2}\in(1,4/3)\), divide its displacement by
\(\sqrt2g\). Its coordinates are one through fifteen,
\(2-u\) at sixteen, and \(u/2\) at seventeen and eighteen.
Its product with \(x\) is at most \(S_{17}\) by ordering.
The actual pressure is at least
\(\sqrt{17-u/2}-\sqrt{(7u-4-3u^2)/2}\).
The first radical exceeds \(4.04\) and the second is at most
\(1/\sqrt{24}<.2042\); hence \(S_{17}>34/9\).
All these finite cases contradict the same prefix rule:
the gains at \(G=34/9\) exceed \(.065,.047\).

For every dotted \(g>1\), put \(u=g^{-2}\).
The actual future reciprocal core \(c+n_{17}/g\)
has norm \(-(1-u)\). Its displacement is
\(U=\sqrt2gP_{17}-n_{17}/g\).
For \(r=17-K\), choose
\[
 Y=U-\sqrt2gP_K,\qquad
 \nu=d_K/(\sqrt2g),\qquad
 Y^2=2r/u-2+u.
\]
The residual is supported after \(K\), so is orthogonal
to the entire already paid reference. Its price is
\(\nu\sqrt2g=d_K\), allowing equality at the terminal.
The gain is exactly
\[
 G_K(u)=\sqrt2d_K\!
 \left[\sqrt{34-u}-\sqrt{2K}+.58-\sqrt{u(1-u)}\right]
 -d_K^2(r-u+u^2/2).
\]
The chord \(\sqrt{34-u}\geq\sqrt{34}-.0864u\) follows from
\(\sqrt{34}>5.8309,\sqrt{33}>5.7445\).
The circle square gives
\(\sqrt{u(1-u)}\leq1/4+u-u^2\).
Using \(\sqrt{28}<5.2916,\sqrt{30}<5.4773\),
the bracket exceeds \(u^2-1.0864u+.8693\) or
\(u^2-1.0864u+.6836\), both positive.
Also \(\sqrt2d_{14}>.19247,\sqrt2d_{15}>.18568\).
Use the upper \(d_K\) only in the negative constant and
the lower \(d_K\) in \(d_K^2(u-u^2/2)\).
This gives
\[
\begin{split}
G_{14}(u)&>.183208395u^2-.190576198u+.111662851,\\
G_{15}(u)&>.177060155u^2-.184483062u+.092398928.
\end{split}
\]
For \(F=au^2-bu+c\), the identity
\(4a(F-\tau)=(2au-b)^2+4a(c-\tau)-b^2\)
has positive last constants
\[
 \frac{808151227201376}{10^{18}}\quad(\tau=.061),\qquad
 \frac{244086609159516}{10^{18}}\quad(\tau=.044).
\]
Thus the full dotted interval gives gains \(>.061,>.044\),
and totals \(>1.000550975,>1.00078547\).
No isolated \(17\)-opposite survives.


\subsection{An exhaustive terminal inventory}
\label{app:B18inventory}
Deleting old \(17\) retains actual \(A_{16}\) on
\(n_1,\ldots,n_{16}\) and independent short \(A_1\) on \(n_{18}\).
The actual finite opposite extends this whole white Gram positively.
A one-node finite enlargement of \(A_{16}\) is \(A_{17},B_{17}\)
or \(D_{17}\); exceptional types have insufficient rank.
Thus the complete white list is
\[
\begin{array}{c|l}
\text{left contact}&\text{short contact}\\ \hline
\text{none}&\text{none or one finite label }m\geq3\\
3\text{ at }j=1,16&\text{none, }3,\text{ or }4\\
4\text{ at }j=1,16&\text{none}\\
3\text{ at }j=2,15&\text{none}.
\end{array}
\]
Two left contacts create a cycle. Joining the short root to a
left-three endpoint gives \(A_{18}\) or \(B_{18}\); joining it
to a left-four or a next-node branch is not finite.
White parallel/dotted contacts are impossible at this finite endpoint.
Own contacts still range over zero, every finite label, parallel
and dotted.

Write a cap as
\[
 c=A p-g\omega_{17}-\eta\omega_j-\delta\omega_{18},
\quad
 \sigma=g+\eta j/17+\delta/\sqrt2,\quad
 \Delta=1-\frac{2\eta^2j(17-j)}{17}-\delta^2>0.
\]
Absent contacts have coefficient zero.
Direct retained-block elimination gives
\[
 A^2=34\sigma^2-\Delta,\qquad
 C_{17}=\sqrt{34}\,\sigma/\sqrt\Delta.
\]
There is no residual cross between the actual left and short blocks.
Own parallel is excluded everywhere by the actual positive null
\(c+n_{17}\). Isolation was excluded above.

For left-four alone, the far and near lengths are
\(17\sqrt2g+1\) and \(17\sqrt2g+16\).
Far own zero has zero time; every positive far own gives length
greater than thirteen, while the near length is already at least
sixteen. For left next-three, the lengths are \(17g+1\) at \(j=2\)
and \(17g+15/2\) at \(j=15\). Far own zero has zero time;
far own three is the actual simple cycle
\(c,n_2,\ldots,n_{17}\). Every larger positive far own exceeds
thirteen. Near positive own gives length at least sixteen;
near own zero has the affine positive kernel with coefficients
one on \(c,n_{14}\), two on \(n_{15},n_{16},n_{17}\),
and \(\sqrt2\) on \(n_{18}\).

Left-three plus short-four has
\(\Delta=1/34\) and \(C_{17}=34g+j+17\geq18\).
Left-three plus short-three has \(\Delta=19/68\).
Own three closes a simple cycle; every larger positive own has
\(g\geq1/\sqrt2,\sigma>1\) and length greater than eleven.
At own zero, the far bridge contains the actual affine path
\((n_{15},n_{16},n_{17},n_{18},c)\), with positive kernel
\((1,2,3,2\sqrt2,\sqrt2)\).
For the near bridge, \(v=(\omega_{16}+\omega_{18})/2\).
Deleting actual short \(18\) gives positive finite \(D_{18}\).
Positive projection checks the released wall and every unknown
wall, so this is the actual short-edge endpoint. Its exact data are
\[
 A^2=23/2+8\sqrt2,\quad \sum v_i=9+8\sqrt2,\quad
 (\sum v_i)^2-18A^2=2,\quad
 C_{18}=8+9\sqrt2/2>14,
\]
contradicting \eqref{app:eq:B18sumbounds}.

For pure far-three at \(j=1\), own zero has negative time,
own three gives its old simple cycle, own four has null
\(c+\sqrt2n_{17}+n_{18}\), and own six has null
\(g c+n_{17}+n_{16}/2\). Parallel was already null.
The remaining owns are five, finite at least seven, or dotted, with
\[
 C_{17}=\frac{\sqrt2}{3}(17g+1/2),\qquad g<4/3.
\]
At \(g\geq4/3\) this would be at least
\(139\sqrt2/18>10.9\).
For pure near-three at \(j=16\),
\(C_{17}=\sqrt2(17g+8)/3\).
Own zero has the positive kernel with coefficients one on
\(c,n_{15}\), two on \(n_{16},n_{17}\), and \(\sqrt2\) on \(n_{18}\).
Own three gives a triangle; four and parallel have the nulls above.
Owns at least seven or dotted have \(g>.9\) and length
\(>10.98362\). Only owns five and six need remain.

Near own five is impossible independently of any head.
For \(g=\phi/2\), the positive actual right-\(B_2\) core is
\[
 \xi=c+2gn_{17}+\sqrt2gn_{18}
     =A p-\sqrt2(g+1/2)P_{16},
\quad A^2=(49\phi+31)/2,\quad \xi^2=1-2g^2<0.
\]
Its cap product is this same negative square and its selected
old products vanish. All other original wall products are
nonpositive; positive time makes it future. The pressure gives
\[
 S_{16}>
 \frac{A-\sqrt{2g^2-1}}{\sqrt2(g+1/2)}>3.7.
\]
Indeed \(A>7.42,\sqrt{2g^2-1}<.556\) and
\(\sqrt2(g+1/2)<1.852\), while
\(7.42-.556>3.7(1.852)\).
The uniform prefix rule gives increments
\(.06306872,.04497024\), contradicting its norm budget.
For near own six, \(g=\sqrt3/2\) gives \(C_{17}>10.5\), hence
\(S_{17}^2>17(19/23)>f_{17}^2\).
Thus this surviving role has strong \(17\) and the isolated
native-five head already proved above.

Finally suppose there is no left contact and
\(\delta=\cos(\pi/m)\) is the short-white magnitude.
The length bound implies
\[
 \frac{g+\delta/\sqrt2}{\sqrt{1-\delta^2}}<1.866 .
\]
For \(m\geq9\), \(\pi^2<10\) and
\(\cos t>1-t^2/2\) give \(\delta>76/81>15/16\);
even at \(g=0\) the ratio square exceeds
\(225/62>1.866^2\). The entire conservative own list is
\[
\begin{array}{c|l}
m&\text{own }g\\ \hline
3&0,\text{ every finite positive own, or }1<g<1.264\\
4&\text{finite own }3,4,5\\
5&0\text{ or finite own }3\\
6,7,8&0 .
\end{array}
\]
For short three,
\(g<1.866(.8665)-.3533<1.264\).
Short four gives \(g<.82\); its own zero has positive null
\(c+n_{17}+\sqrt2n_{18}\).
Short five gives \(g<.53\) from
\(\delta>.809,\sqrt{1-\delta^2}<.588,\delta/\sqrt2>.571\).
For short at least six, \(g<.933-.6<1/2\).
Every null above is a positive actual-root combination with
positive cap time, zero selected products and nonpositive product
against every remaining original wall. It is therefore a nonzero
feasible future ray, forbidden by compactness.

This proves the exhaustive partition into a far-three helper,
near-three/own-six, or the displayed no-left short-contact family.

\subsection{Finite far-helper owns from raw, unspent references}
\label{app:B18finitefar}
The far cap is
\[
 c=A p-\sqrt2gP_{17}-e_1/\sqrt2,\qquad
 A^2=34g^2+2g-\tfrac12.
\]
Use the three raw contracts of Section~\ref{app:B18references},
before any head graft. Their terminal prices are at least
\(d_K\), and their full first coordinates are bounded there.
Pay \(\mu_{14}=.136,\mu_{15}=.1313\) for a new pressure
\(\ip xV>G\), with \(Y=V-P_K\).
The full square retains \(\ip wY\); no absent \(P_{14}\)
coefficient is borrowed when the actual strong knot is \(15\).

For own five, the actual \(H_4\) correction
\[
 \eta=c+(3g/2)n_{17}+gn_{16}+(g/2)n_{15}=A p-U
\]
has positive square \(\delta=(5-3\phi)/8\).
Use the actual finite point \(p+(A/\delta)\eta\).
Its selected products vanish, its correction coefficients are
positive and every unselected original wall remains nonpositive.
Its pressure exceeds \(A-\sqrt\delta\).
For \(b=1/(2g)=1/\phi<.619\), the vector \(V=U/(\sqrt2g)\)
has coordinates \(1+b\) at one, one through fourteen, and
\(3/4\) at fifteen through eighteen. Its pressure exceeds \(4.10\):
\[
 A^2=(19\phi+16)/2,\quad
 A>4.834,\quad \sqrt\delta<.136,\quad \sqrt2g<1.145.
\]
The complete residual data are
\[
\begin{array}{c|cc}
K&Y^2<&\ip wY\\ \hline
14&2.634&bw_1\\
15&2.134&bw_1-w_{15}/4 .
\end{array}
\]
Use \(w_{15}>.1313\) and the actual raw row's \(L,w_1\).
The resulting exact strict norm lower bounds are
\[
\begin{array}{c|cc}
\text{initial five}&K=14&K=15\\ \hline
\text{strong}&1.0217895424&1.02014330566\\
\text{helper}&1.020882792&1.01938584884\\
\text{isolated dotted weak}&1.0258028&1.02332542274 .
\end{array}
\]
For example the first column is obtained from
\(L+.272(4.10-3.332-.619w_1)-.018496(2.634)\).
Every row is the same original square with its entire first-coordinate
cross retained.

For every finite own at least seven, put \(u=g^{-2}\in(1,4/3)\).
The positive actual core
\[
 \xi=c+g^{-1}n_{17}+2(g^{-1}-g)n_{16}=A p-U
\]
has norm \(-(u-1)(4-3u)/u\), zero cap/own products
and negative \(n_{16}\) product. Its unknown wall signs and
future sheet follow from positive actual coefficients and \(A>0\).
For \(b=1/(2g)\in(.5,.556)\), \(V=U/(\sqrt2g)\)
has coordinates
\[
 1+b,\quad 1\ (2,\ldots,15),\quad 2-u,\quad u/2,\ u/2 .
\]
Its pressure is
\[
 G=\sqrt{17+2b-b^2}-\sqrt{(7u-4-3u^2)/2}
   >\sqrt{71/4}-1/\sqrt{24}>4.0088 .
\]
The residuals satisfy \(Y_{15}^2<7/4,\ Y_{14}^2<11/4\);
both full crosses are \(bw_1\).
Indeed
\(Y_{15}^2=4-15u/4+3u^2/2\), whose difference from \(7/4\)
is \((3/2)(u-1)(u-3/2)<0\).
The three initial-five rows give
\[
\begin{array}{c|cc}
 &K=14&K=15\\ \hline
\text{strong}&1.0010956736&1.00023617738\\
\text{helper}&1.003278544&1.0024615677\\
\text{isolated dotted weak}&1.005311136&1.0036135113 .
\end{array}
\]
The tight entry uses precisely
\[
 .9403998+.2626(4.0088-3.463-.556(.3652))
       -.01723969(7/4)=1.00023617738.
\]
The weaker energy \(.94006\) or a terminal weight \(.13\)
would not justify it. All finite far-helper owns are excluded
without a head or short-opposite assumption.

\subsection{A single table for every actual short opposite}
\label{app:B18shorttable}
Deleting actual \(18\) retains \(A_{17}\).
The finite one-node classification gives precisely isolation, an
endpoint bond three or four at \(j=1,17\), or a next-to-end bond
three at \(j=2,16\). A second contact creates a cycle; the other
branching trees and exceptional diagrams have insufficient rank.
For the actual short opposite \(d=Bp-v\), write
\(\delta=-\ip d{n_{18}}\geq0\). At this stage its own contact
can be zero, any finite label, parallel, or dotted. The complete data are
\[
\begin{array}{c|c|c|c}
\text{white role}&v&B^2&C_{18}\\ \hline
\text{isolated}&\delta P_{18}&18\delta^2-1&3\sqrt2\delta\\
\text{end three }j&\delta P_{18}+P_j/\sqrt2&
18\delta^2+\sqrt2j\delta+j/2-1&
(18\sqrt2\delta+j)/\sqrt{19}\\
\text{end four }j&\delta P_{18}+P_j&
18\delta^2+2j\delta+j-1&18\delta+j\\
\text{next three }j&\delta P_{18}+P_j/\sqrt2&
18\delta^2+\sqrt2j\delta+j/2-1&9\sqrt2\delta+j/2 .
\end{array}
\]
For example the old short residual is \(1/18\), and clearing a
white three or four leaves residual
\(1-j(18-j)/36\) or \(1-j(18-j)/18\), respectively.
The edge formula gives the last column; thus the table includes
the actual short normalization.

The unconditional bound \(C_{18}<9.81\) removes near end-four
and every positive-own far end-four. The latter has \(B=0\)
at own zero. Near next-three has \(C_{18}=9\sqrt2\delta+8\),
so a positive own is already too long; at own zero its null is
\[
 d+n_{15}+2n_{16}+2n_{17}+\sqrt2n_{18}.
\]
Far next-three has \(B=0\) at own zero and \(\delta<.7\)
by length, so its only possible positive native own is three.
Its null is
\[
 \sqrt2d+2\sqrt2n_{18}+3n_{17}+2n_{16}+n_{15}.
\]
All coefficients are positive, the cap time is positive,
selected products vanish and every other original-wall product
is nonpositive termwise. These are forbidden future feasible nulls.

We shall apply the remaining three rows twice. In both applications
the selected seventeenth cap clears positively to
\[
 \bar c=Ap-(Q/17)P_{17},\quad
 a=A^2,\quad Q^2=17a+9,\quad \bar c^2=r=9/17,\quad a<63.
\]
For near end-three the actual mutual magnitude is
\(\lambda=AB-Q(\delta+1/\sqrt2)\geq0\).
The length bound gives \(\delta<1.1\), whereas
\[
 aB^2-Q^2(\delta+1/\sqrt2)^2
 =(a-9)\delta^2-9\sqrt2\delta-a-\tfrac92
 <.21a-15.39<0.
\]
This contradicts acuteness, including own zero and every allowed
dotted value. Only far end-three and isolation remain.

\subsection{Dotted far helpers: the complete mutual split}
\label{app:B18dottedfar}
Let
\[
 c=Ap-\sqrt2gP_{17}-e_1/\sqrt2,\quad
 1<g<4/3,\quad a=34g^2+2g-\tfrac12,\quad
 Q=(34g+1)/\sqrt2 .
\]
Here \(71/2<a<1127/18\), and the positive actual correction is
\(\bar c=c+\sum_{i=1}^{16}(17-i)n_i/17\).
Put \(G_{17}=3.760,\ G_{18}=3.830\), and
\[
 t=Q/A=\sqrt{17+9/a},\quad K=A\sqrt{18}-Q,\quad
 H=A-G_{17}Q/17,\quad
 J_0=(t-G_{18})/K.
\]
Then \(4.14<t<4.18,\ K,H,J_0>0\), and
\begin{equation}\label{eq:B18quadraticframe}
 4HJ_0>1,\qquad H+J_0y^2-y>0\quad(y\in\mathbb R).
\end{equation}
Indeed the first inequality is equivalent to positivity of
\(F(t)=4(t-3.83)(1-94t/425)+t-\sqrt{18}\).
Replace \(\sqrt{18}\) by its strict upper bound \(4.243\).
The new polynomial has value \(3333/2125000>0\) at \(4.14\);
its difference from that value is
\[
 (t-4.14)\{5-(376/425)(t+4.14-3.83)\}>0,
\]
since the brace exceeds \(10919/10625\). The second inequality
is the completed square
\((2J_0y-1)^2+4HJ_0-1>0\).

We also record a useful full-projection test.
Suppose an inverse-positive full actual projection has ray
\((1+T)p-LP_{17}-Me_{18}\), with \(T>200\) and \(L>M>0\).
Write \(k=M/L\). Its norm and the original farthest comparison give
\[
 S_{17}+kx_{18}\geq
 \sqrt{17+k^2}\frac{\sqrt T}{\sqrt{1+T}+1}
 >\sqrt{221/15}>3.83.
\]
For the strict bound use \(\sqrt{1+T}>14\), so the squared
last ratio is \(>13/15\). The left side is
\((1-k)S_{17}+kS_{18}<3.83\), a contradiction.
The existence, future sheet and unknown-wall signs of this point
follow from its full positive actual projection.

First suppose that the short opposite is far end-three:
\[
 d=Bp-\delta P_{18}-e_1/\sqrt2,\qquad
 B^2=18\delta^2+\sqrt2\delta-\tfrac12 .
\]
Own zero has negative time. Own three has the affine \(F_4\)
null above, and owns four, six and parallel have respectively
\[
 d+\sqrt2n_{18}+n_{17},\qquad
 d+(\sqrt3/2)n_{18}+n_1/2,\qquad d+n_{18}.
\]
These positive actual nulls remain feasible against every original wall.
For every remaining own put \(\lambda=-\ip cd\geq0\) and
\(v_0=1/(17\sqrt2)\). Positive clearing gives
\[
 \bar d=Bp-(\delta+v_0)P_{17}-\delta e_{18},\quad
 \bar d^2=r,\quad
 \rho=-\ip{\bar c}{\bar d}=\lambda+8/17
       =AB-Q(\delta+v_0).
\]
The offset here is \(8/17\). If \(\lambda>0\), the actual native
gap gives \(\lambda\geq1/2\), hence \(\rho>r\).
The positive actual core \(\Xi=\bar c+(\rho/r)\bar d\)
has square \(r-\rho^2/r<0\), positive time, zero \(d\) product,
negative \(c\) product, and zero old \(A_{16}\) products.
All other products are nonpositive by acuteness.
Its pressure and the two original sum bounds would require
\[
 H+(\rho/r)J-\sqrt{\rho^2/r-r}\leq0,\qquad
 J=B-G_{17}v_0-G_{18}\delta .
\]
But \(J=\rho/A+(t-G_{18})\delta+(t-G_{17})v_0
>(t-G_{18})\delta\).
Positive squaring gives \(B<\sqrt{18}\delta+1/6\), and
\[
 36(g+1/34)^2-a
 =2g^2+2g/17+\tfrac12+9/289>0 .
\]
As \(Qv_0=g+1/34\), this implies \(\rho<K\delta\)
and \(J>J_0\rho\).
With \(y=\sqrt{\rho^2/r-r}\), the required gap exceeds
\(H+J_0(y^2+r)-y>0\) by \eqref{eq:B18quadraticframe}.

If \(\lambda=0\), the cleared two-cap Gram has diagonal \(9/17\),
cross \(-8/17\), and inverse
\(\left(\begin{smallmatrix}9&8\\8&9\end{smallmatrix}\right)\).
Its positive projection has
\[
 \alpha=9A+8B,\quad \chi=8A+9B,\quad
 T=9a+16AB+9B^2>319.5,\quad
 L=\alpha Q/17+\chi(\delta+v_0)>M=\chi\delta>0 .
\]
The full-projection test excludes it. This handles every far-three
short own and every unbounded actual mutual range.

Now let \(d=Bp-\delta P_{18}\) be isolated.
Acuteness \(AB-Q\delta\geq0\) implies
\(\delta^2\geq a/(a-9)>1\), so its own is derived to be dotted.
Put \(\lambda=AB-Q\delta\).
The cleared Gram is
\(\left(\begin{smallmatrix}r&-\lambda\\-\lambda&1\end{smallmatrix}\right)\).
At \(\lambda^2=r\), \(\bar c+\lambda d\) is a forbidden positive
future feasible null. For \(\lambda^2>r\) the same core is negative.
Writing \(J=B-G_{18}\delta\), acuteness and \(B<\sqrt{18}\delta\)
give \(J>(t-G_{18})\delta,\ \lambda<K\delta\).
Its required pressure gap is
\(H+\lambda J-\sqrt{\lambda^2-r}\), contradicted by
\eqref{eq:B18quadraticframe}.
If \(0<\lambda^2<r\), native separation gives \(\lambda\geq1/2\).
The positive full projection has
\[
 T=\frac{a+2\lambda AB+rB^2}{r-\lambda^2}
 >\frac{68.5}{19/68}=\frac{4658}{19}>200 ,
\]
using \(B^2>17,\ AB>24\).
Its coefficients
\(\alpha=(A+\lambda B)/(r-\lambda^2)\),
\(\chi=(\lambda A+rB)/(r-\lambda^2)\)
are positive, with \(L=\alpha Q/17+\chi\delta>M=\chi\delta>0\).
The full-projection test excludes this entire interval too.
Only \(\lambda=0\) remains.

\subsection{The orthogonal neighbor and its quantitative comparison}
\label{app:B18A17comparison}
For that last mutual value,
\[
 \delta^2=\frac a{a-9},\quad B=\frac Q{\sqrt{a-9}},\quad w=9/a,\quad
 q=\frac{p-(A/Q)(P_{17}+we_{18})}{\sqrt D},\qquad
 D=\frac{9(a-9)}{aQ^2}.
\]
The positive actual projection has full forest
\(A_{17}(c,n_1,\ldots,n_{16})+I(d)\).
The other endpoint retaining its whole \(A_{17}\) is the positive
finite projection on \(A_{16}+c+n_{18}\). These endpoints are distinct
because \(\ip d{n_{18}}=-\delta<-1\).
Consequently \(n_{18}\) is the actual opposite of \(d\), orthogonal
to the entire new \(A_{17}\), with
\[
 h_I=\delta^2-1=\frac9{a-9}<18/53<1/2.
\]
All these are full-dimensional points satisfying every original wall;
no face heredity is used.

Write the same original \(z=E_q(q-y-\eta d)\), where \(y\)
is in the \(A_{17}\) space. We prove the following three bounds,
separately for old \(17\) strong and weak:
\begin{equation}\label{eq:B18A17interface}
 E/E_q<1.061,\qquad \eta^2>.1164,\qquad
 \text{a signed \(A_{17}\) endpoint coordinate of \(y\) exceeds } .51 .
\end{equation}
Every actual chamber point, including each auxiliary point used below,
is then compared at \(q\) with factor \(<1.061\), because its height
relative to the original \(z\) is at most \(E\).
The point \(q\) is never made farthest.

Suppose first that \(S_{17}>f_{17}=s\), so the earlier head argument
gives the actual isolated native-five head and strong \(5,6,7\).
The original reference through strong \(17\), with all internal
weak knots omitted, has energy \(>.97253117\),
first-five mean \(b<.3652\), fifth capacity \(>.151\), and
terminal capacity \(b_0=\sqrt{17}-4\in(.1231,.1232)\).
The actual head foot is \(Z_0=P_4/4\), with pressure
\(h_0>.39307\). Put \(Y_0=Z_0-P_5/5\), so \(Y_0^2=1/20\).
Paying \(.56\) costs \(.112<.151\) and gains \(>.0155344\).
The paid effective energy is \(L>.988\). Retaining the terminal
surplus in this same square gives
\[
 2b_0(S_{17}-s)+x_{18}^2<.012,\qquad
 S_{18}<s+.012/(2b_0)+b_0/2<3.82335 .
\]

We next show \(g>1.12\). For \(1<g\leq1.12\) the reciprocal core
\(c+n_{17}/g\) has pressure \(A-\sqrt{1-g^{-2}}\).
Its residual after payment from \(P_{17}\) and \(P_5\) has square
\(.4+g^{-2}\), and cross \(1/(20\sqrt2)\) with \(Y_0\).
The pressure-minus-reference deficit is
\[
 D_H=A-\sqrt{1-g^{-2}}-g(\sqrt{34}-.58)
       -(\sqrt{10}/10-.058)+b_0/(g\sqrt2).
\]
It decreases on this interval: \(A'^2=34+18/a<5.875^2\);
the radical derivative exceeds \(1.5\), and
\(\sqrt{34}-.58>5.25\), so \(D_H'<-.875\).
At \(g=1.1\) use \(A>6.5452,\ \sqrt{1-g^{-2}}<.4166\),
the third term \(<5.7761\), and the last term \(>.079\);
at \(1.12\) use \(6.66255,.45035,5.88112,.07767\).
Also \(\sqrt{10}/10-.058<.258228\).
Thus \(D_H>.173\) on \((1,1.1]\), and \(>.150\) on
\([1.1,1.12]\).
With residual coefficient \(\nu=.075,.077\) respectively, the
full gain over \(L\) is
\[
 2\nu D_H-\nu^2(.4+g^{-2})-.56\nu/(10\sqrt2)
 >.015075,\quad >.0127780815 .
\]
Both contradict \(L>.988\). The individual terminal price
\(\nu\sqrt2g<.1231\) fits \(b_0\), while
\(.112+\nu/(5\sqrt2)<.123<.151\) fits the fifth anchor.
Hence \(a>44.3896\), and \(k=9/a<.20276\).

Suppose \(E_q/E\leq20/21\).
The explicit \(q\) formula gives
\[
 S_{17}+kx_{18}\geq
 R(k):=\sqrt{17+k}-(20/21)\sqrt{k(1-k)}>3.764 .
\]
Indeed \(R\) decreases on \(0<k<.20276\): its positive derivative
term is \(<.122\) and the magnitude of the negative term is
\(>(20/21).594\). At the upper endpoint use
\(\sqrt{17.20276}>4.1476,\ \sqrt{.20276(1-.20276)}<.4021\).
The independent strong \(K=11,12\) references through strong \(5,6,7\),
after every permitted hole, have raw energies
\(>.86375298,>.88550238\), terminal capacities \(>.1543,>.1474\),
and fifth capacity \(>.151\).
After the same \(.56\) head payment, pay \(.14\) from that terminal
toward \(P_{17}+ke_{18}\).
Its residual is orthogonal to the head residual and to the reference,
of square \(17+k^2-K\), with \(k^2<.0412\). Since
\(f_{11}<2.907,\ f_{12}<3.055\), the full original squares give
\[
\begin{split}
 .86375298+.0155344+.28(3.764-2.907)-.0196(6.0412)
 &=1.00083986,\\
 .88550238+.0155344+.28(3.764-3.055)-.0196(5.0412)
 &=1.00074926 .
\end{split}
\]
Thus \(E_q/E>20/21\), so \(E/E_q<1.05\).
Moreover \(a<64,\ \delta>8/\sqrt{55}>1.078,\ Q/A>4.14\) imply
\[
 \eta=(E/E_q)\delta(Q/A-S_{18})
 >1.078(4.14-3.82335)>.3413,\qquad \eta^2>.1164 .
\]
The known actual \(q\)-edge releasing \(c\) retains \(A_{16}+d\);
its other endpoint is \(A_{16}+n_{17}+d\). Clearing that block
leaves residuals \(9/17\) and cross \(-g-1/34\), so its cosine is
\(C'=(34g+1)/18>977/450>2.17\).
The signed endpoint coordinate is therefore greater than
\[
 \sqrt{17/18}\frac{2.17-1.05}{\sqrt{2.17^2-1}}
 >.9718(1.12)/1.926>.56 .
\]
This proves \eqref{eq:B18A17interface} in the strong case.

If old \(17\) is weak, \(S_{17}\leq f_{17}<3.713\), and
\[
 E_q/E=\frac{\sqrt{17+w}-S_{17}-wx_{18}}{\sqrt{w(1-w)}}.
\]
The two original bounds give \(S_{17}+wx_{18}<3.713+.117w\).
If \(g\geq1.1\), then \(w\leq25/119<.211\).
On this range \(\sqrt{17+w}>4.123+.12w\), since
\(\sqrt{17+w}+\sqrt{17}<25/3\), while
\(\sqrt{w(1-w)}<.408\). The displayed ratio exceeds
\((.410+.003w)/.408>1\), contrary to \(E_q\leq E\).
Thus \(1<g<1.1\) and \(25/119<w<18/71\).
The same chord holds, \(w>.21\), and the denominator is
\(<\sqrt{954}/71<.4351\).
Consequently
\[
 E_q/E>.41063/.4351>.943,\qquad E/E_q<1000/943<1.061.
\]
Here \(\delta^2>119/94>81/64\) and
\(\sqrt{17+w}>4.1485\), so
\(\eta>(9/8)(4.1485-3.830)=.3583125\), hence \(\eta^2>.128\).
Finally \(C'>35/18>1.94\), and its endpoint coordinate exceeds
\[
 .9718(1.94-1.061)/1.663>.51 .
\]
All three conclusions of \eqref{eq:B18A17interface} follow.

\subsection{One local \(A_{17}+I\) contradiction for both height branches}
\label{app:B18localA17}
We prove the common consumer of \eqref{eq:B18A17interface}.
The hypotheses are an actual whole \(A_{17}+I\) vertex \(q\),
comparison factor \(\kappa<1.061\), a known isolated actual \(I\)
opposite with parameter \(h_I<1/2\), and
\[
 \norm y^2+\eta^2<1,\qquad \eta^2>.1164,\qquad
 \text{one signed endpoint coordinate of \(y\) is }>.51 .
\]
We use the finite-white geometry proved in
Appendices~\ref{app:Acentralwindow} and~\ref{app:Aendpointcores},
not their farthest-vertex energy conclusions.
Here are the geometry and the changed numerical comparisons explicitly.

Realize \(A_{17}\) in the zero-sum subspace of \(\mathbb R^{18}\).
For its unit roots \(m_i\) and duals \(\Omega_i\),
\[
 D_i=i(18-i)/9,\quad C_{ij}=i(18-j)/9\ (i\leq j),\quad
 h_i=\ip y{\Omega_i},\quad
 e_A=\tfrac12\sum_{i=0}^{17}(h_{i+1}-h_i)^2 .
\]
The sequence \(h\) is strictly concave.
A site \(5\leq i\leq13\) is weak if \(h_i\leq\sqrt{D_i}-.65\).
Deleting such \(i\) retains two paths of ranks at least four and
the entire actual \(I\). The finite one-node list is isolation;
one simple bond at an end, \(D\) or permitted \(E\) position;
an end-four; or two end-threes, one on each path.
A joined component has rank seventeen and is classical.
The retained foreign \(I\) cost \(Q_I\) is less than one and,
if nonzero, at least \(1/4\).

Own zero is removed geometrically before any comparison:
internal-three roles have affine arms \(1,2,5\) or \(1,3,3\);
internal-four gives the \(4,3,3\) star; two whites give an affine
cycle. A cap without an \(A\) contact cannot close the edge.
End-four requires \(Q_I<1/m\leq1/5\), hence \(Q_I=0\) and
negative time square. In the remaining end-three/node-two cases
the full old-component norms give \(A_0^2<Q_I\).
The known actual \(I\) opposite is
\(\sqrt{h_I}q-\sqrt{1+h_I}d\); acuteness would require
\(A_0\sqrt{h_I}\geq\sqrt{Q_I}\sqrt{1+h_I}\), impossible.
These are the zero-own calculations of
Appendix~\ref{app:Acentralwindow}, with this actual opposite
replacing its farthest-only orthogonal-opposite step.

For a white coefficient vector \(b\), all remaining length data
are computed by
\[
 \Delta=1-b^\mathsf TG_{\rm retained}^{-1}b-Q_I,\qquad
 \sigma=\gamma+
  \ip{\Omega_i}{\sum_{j\ne i}b_j\Omega_j}/D_i,\qquad
 C=\sqrt{D_i}\,\sigma/\sqrt\Delta .
\]
The finite-tree calculation in that appendix gives ratio at least
\(3/2\) for near-three/own at least five, \(D/E\), end-four and
two-white rows, after their actual affine small-own boundaries.
Positive foreign cost increases this ratio.
Foreign far-three has ratio \(>227/(25\sqrt{35})>3/2\);
foreign near-four has squared ratio \(>121/35\).
A near-four with side order at least eight has ratio \(>761/500\),
and side order seven has ratio \(>3/2\).
For a pure helper with side order \(m\),
\[
 \sigma=\gamma+1/(2m),\quad \Delta=(m+1)/(2m),\quad
 \sigma^2/\Delta-(4\gamma-2)
 =(\gamma-1-1/(2m))^2/\Delta\geq0 .
\]
In particular \(\gamma\geq17/16\) also gives ratio at least \(3/2\).

These geometric bounds are sufficient under the new comparison.
For \(u=\sqrt{D_i}\in(\sqrt{65}/3,3]\), the positive-square test at
\(C=3u/2\) is
\[
 (3u/2-1.061)^2-(1-.65/u)^2(9u^2/4-1)
 =F(u)=1.175096-.258u-1.3/u+.4225/u^2>0 .
\]
Indeed \(F(3)>.0147\), and
\[
 F(u)-F(3)=(3-u)
 \{.258-1.3/(3u)+.4225(3+u)/(9u^2)\}\geq0,
\]
as \(u>8/3\). The actual directional floor
\((C-\kappa)/\sqrt{C^2-1}\) increases here because \(\kappa C>1\).
Every listed ordinary role is therefore strong.

For isolated \(A\) contacts, after positive clearing of the whole
\(I\), set \(U=(1-Q_I)/\gamma^2\);
\(A_0^2=\gamma^2(D_i-U)\).
For \(0<U<1\) the positive reciprocal core gives
\[
 h_i>\sqrt{D_i-U}-1.061\sqrt{U(1-U)}.
\]
Its loss is less than
\(.194U+1.061\sqrt{U(1-U)}
 <(.194+1.079)/2=.6365<.65\).
For \(U>1\) the positive finite projection gives loss
\(<.194+1.061\sqrt{U-1}\).
For foreign \(Q_I>0\), \(U-1\leq5-3\phi<(5/12)^2\);
for pure finite own at least eight, the same bound follows from
\(\tan(\pi/8)<5/12\). The loss is \(<.637\).
The \(U=1\) core is a positive future feasible null.
Pure own-five uses the actual \(H_4\) finite point:
its old loss is \(<25/76+5/24\), and replacing comparison one by
\(1.061\) adds
\[
 .061\sqrt{\delta_5}\sqrt{1+\delta_5/A_0^2}/\gamma<.016,
 \quad \delta_5=(7-3\sqrt5)/16<1/36,\quad
 \gamma>.8,\quad A_0^2>3.6 .
\]
Its total loss is \(<.554\).
Pure seven has old loss \(<2500/10368+.36\) and added loss
\(<.061(.36)=.02196\), hence total \(<.624\).
Own three, four, six and parallel have the actual nulls proved in
the same appendix. Thus no isolated or foreign role remains weak.

The provisional weak list is exactly pure near-four at \(5,6\)
or reflection, or a pure far helper with
\[
 v=\gamma\Omega_i+\Omega_e/2,\quad e=1\text{ or }17,\quad
 \phi/2\leq\gamma<17/16,\quad Q_I=0 .
\]
This cutoff concerns only weak helpers.
Their nonadjacency is the original-vector calculation of
Appendix~\ref{app:Acentralwindow}. To specify its complete domain,
for adjacent sites and owns \(\alpha,\beta\in(.809,1.08)\)
the common/outward/reversed products have lower bounds
\[
 1-1.08^2,\qquad
 \tfrac12-\alpha\beta+(\alpha-1)(\beta-1)/18,\qquad
 \tfrac12-\alpha\beta+(\alpha+1)(\beta+1)/18 .
\]
Native gaps leave common mutual zero or opposite mutual zero/simple.
The reversed zero case has numerator
\[
 18F=4\ell s\alpha^2\beta^2-(\ell+1)(s+1)(\alpha^2+\beta^2)
       -2\alpha\beta-(\ell+1)\alpha-(s+1)\beta+5,
 \quad \ell+s=17,\quad\ell,s\geq5 .
\]
It increases in both owns on the permitted range. Unless both are
five, one is at least seven and exceeds \(.9005\);
concavity in \(\ell\) reduces to endpoints, whose smaller value is
\(.05526706486>0\) at \((.809,.9005)\).
The both-five case is the actual distance-one determinant obstruction
proved there, on its declared \(N\geq15\) domain.
The outward zero case has determinant at least \(404/225>0\).
The resulting mixed projection is positive with positive inverse;
adding the entire old \(I\) block leaves it positive.
Its actual future finite point is feasible against every original wall.
The complete inverse calculation bounds the weighted loss by
\(\kappa(7/16)+1/10<.5641875\) for a common helper and
\(\kappa(3/10)+1/10<.4183\) for opposite helpers.
Both are less than \(.65\), proving nonadjacency before near-four
has been removed.

A near-four has an actual mate at \(j=4,5\), obtained by the full
positive \(A_{j-1}+B_{18-j}+I\) projection.
Its geometric cosine exceeds ten, so its new demand exceeds
\((10-1.061)/\sqrt{99}\), giving \(h_4>2.24\) or \(h_5>2.41\).
At the other end there is a mate or a strong \(12\) or \(13\)
by nonadjacency, with demand \(>2.178\) or \(>2.037\).
Disjoint end-segment Cauchy has worst case
\[
 e_A>\tfrac12(2.41^2/5+2.178^2/6)=.976117>.8836,
\]
contrary to \(\eta^2>.1164\). The other placements are stronger.
Every remaining weak site is therefore a nonadjacent pure helper.

Here is the entire reference, including every allowed internal hole.
Take symmetric heights
\[
 H_5=2.077,\ H_6=2.218,\ H_7=2.314,\ H_8=2.371,\ H_9=2.389,
 \qquad H_{18-i}=H_i ,
\]
linear from zero to \(5\), and scale them by \(\rho_0=.98\).
Positive squares \((\rho_0H_i+.65)^2<D_i\) put them below the
strong floors. Their positive successive slopes before scaling are
\(.4154,.141,.096,.057,.018\).
Flatten the first block through \(6\) if \(5\) is weak and
reflect if \(13\) is weak; replace each remaining weak knot by
its neighbor chord. Deleting \(j\) costs
\((H_{j-1}-2H_j+H_{j+1})^2/4\) in reference energy.
Nonadjacency makes these edge pairs disjoint. The following bounds
cover all hole sets without enumerating subsets:
\[
\begin{array}{c|c|c|c|c}
\text{weak ends}&\text{unscaled initial energy}&
 \text{holes}&\text{loss per hole}&\text{scaled energy}\\ \hline
0&.8954558&\leq4&\leq.00050625&>.858\\
1&.8640827333\ldots&\leq3&\leq.00050625&>.8284\\
2&.8327096666\ldots&\leq3&\leq.00038025&>.7986 .
\end{array}
\]
The initial values are exactly, with
\(a=2.077,\ b=2.218,\ T=.096^2+.057^2+.018^2\),
\[
 a^2/5+(b-a)^2+T,\quad
 \tfrac12\{b^2/6+a^2/5+(b-a)^2+2T\},\quad b^2/6+T .
\]
The available individual end-anchor capacities are at least
\(\rho_0(821/6000)\) on a six-block or
\(\rho_0(343/2500)\) on a five-block, both \(>.134\).
The signed block means satisfy \(b_5<.28812,\ b_6<.98(.2614)\).
Every remaining positive coefficient is on an actual strong site;
no weak floor or \(I\) energy is in these references.

For completeness we now give the feet, deficits and all mutual
residual terms needed to use these capacities.
At a weak end let \(m_i,m_e,m_j\) be its actual own, far-helper
and helper-neighbor roots. Write
\(U=v-k_0m_i-\ell m_e-mm_j\). The complete own dispatch is
\[
\begin{array}{c|c|c|c}
\text{own}&(k_0,\ell,m)&\text{corrected square}&
 \text{pressure with comparison one}\\ \hline
5&(\gamma,2/3,1/3)&\delta=2/3-\gamma^2>0&
 >A-\sqrt\delta\\
\text{finite }\geq7&(\gamma,2(1-\gamma^2),0)&
 -(1-\gamma^2)(4\gamma^2-3)&
 A-\sqrt{(1-\gamma^2)(4\gamma^2-3)}\\
\text{dotted}&(1/\gamma,0,0)&-(1-\gamma^{-2})&
 A-\sqrt{1-\gamma^{-2}} .
\end{array}
\]
In the first row the point is
\(q+(A/\delta)(b+k_0m_i+\ell m_e+mm_j)\);
the spacelike corrected normal itself is not a point.
The other rows are future negative cores with zero cap product,
nonpositive selected old products and all unknown products
nonpositive termwise. Time coefficients are positive.
All displacements are pure \(A_{17}\).

For a far orientation the time is
\(A_0(\gamma)^2=(65/9)\gamma^2+(5/9)\gamma-19/36\);
near adds \(8\gamma/9\).
On two six-blocks the projections are
\[
 \Pi U=(5\gamma/6)\Omega_6+\Omega_{12}/12\quad\text{(far)},\qquad
 \Pi U=(5\gamma/6+1/12)\Omega_6\quad\text{(near)} .
\]
Let \(Y=U-\Pi U\). Its sum is zero separately on each end block,
so \(\ip wY=0\) after every internal hole.
The far residual squares are exactly
\[
 Q_5=\tfrac23\gamma^2+\tfrac1{12},\quad
 Q_f=4\gamma^4-\tfrac{16}3\gamma^2+\tfrac{29}{12},\quad
 Q_d=\tfrac53\gamma^2-\tfrac{19}{12}+\gamma^{-2};
\]
near adds \(\gamma/3\).
Every finite far square is \(<13/12\), and every dotted far square
is \(<148/125\).

We give rational seals for the pressure deficit against the unscaled
pairing \(2.218(5\gamma/6+1/12)\).
For owns five, seven, eight and at least nine the respective lower bounds
are
\[
\begin{array}{c|c|c|c}
\text{own}&\text{own bracket}&\text{time/radical bounds}&
 \text{deficit lower bound}\\ \hline
5&\gamma<.81&A_0>2.155,\ \sqrt\delta<.111&
 2.155-.111-2.218(91/120)>.362\\
7&.9009<\gamma<.901&A_0>2.4154,\ R<.2158&
 2.4154-.2158-2.218(1001/1200)>.3494\\
8&.9238<\gamma<.9239&A_0>2.4795,\ R<.2464&
 2.4795-.2464-2.218(.85325)>.34059\\
\geq9&\gamma>.93&A_0(.93)>2.497,\ R\leq1/4&
 2.497-.25-2.218(103/120)>.3432 .
\end{array}
\]
All root bounds are positive-square comparisons, using
\(\phi\), the native-seven cubic \(8u^3-4u^2-4u+1\), and
\(\gamma_8^2=(2+\sqrt2)/4\).
For the last row the derivative applies to the lower envelope
\(A_0-\tfrac14-2.218(5\gamma/6+1/12)\):
\(A_0'>\sqrt{65}/3>5(2.218)/6\).
It is not a monotonicity claim for the actual varying-radical deficit.
Thus the finite deficit is \(>.3405\).
For dotted owns split the whole interval \((1,17/16)\) at
\(1.03,837/800,1687/1600\).
The increasing function \(A_0-2.218(5\gamma/6+1/12)\) has
derivative \(>.839\).
Time lower bounds at the first three left endpoints are
\(2.692,2.775,2.82\); radical upper bounds for the four intervals
are \(.24,.295,.317,.338\).
The fourth lower function value is the third plus
\(.839(13/1600)\).
The resulting four deficits exceed
\(.4188,.3913,.3843,.37016\).
This covers the entire dotted range, with deficit \(>.37\);
near times are larger.

Changing comparison from one to \(1.061\) costs \(<.0071\)
for five, at most \(.01525\) for finite at least seven, and
\(<.0207\) for dotted. Indeed the relevant extra terms are
\(.061\sqrt\delta\sqrt{1+\delta/A^2}\),
\(.061/4\), and \(.061(.338)\); for five
\(A>2,\sqrt\delta<.115,\delta<.013\).
Scaling the reference down gains at least
\(.02(2.218)(.7575)=.0336027\).
Thus every six/six deficit is, before rounding, strictly greater than
\[
 .3405-.01525+.0336027=.3588527>.358 .
\]
For one weak end, on six/five blocks the exact far projection is
\[
 \Pi U=(5\gamma/6)\Omega_6+\Omega_{13}/10 .
\]
Relative to six/six, the pairing increases by
\(\rho_0(343/15000)=.022409333\ldots\)
and the residual square decreases by \(1/60\).
The new far deficit is therefore strictly greater than
\[
 .3588527-.98(343/15000)=.336443366\ldots>.336 .
\]
Near is unchanged. In particular the opposite price here is
\(\mu/10\), not \(\mu/12\).

For two feet we retain their entire cross.
Within each length-six block decompose the own and helper residuals as
\[
 a(\gamma,k_0)=\frac{\gamma\sqrt2}{6}(P_5-5e_6)-k_0m_5,\qquad
 b(\ell,m)=\frac{e_1-P_6/6}{\sqrt2}-\ell m_1-mm_2 .
\]
Direct products give
\[
 a^2=5\gamma^2/3+k_0^2-2\gamma k_0,\quad
 b^2=5/12+\ell^2+m^2-\ell m-\ell\leq5/12,\quad
 \ip{a_i}{b_j}=\gamma_i/6 .
\]
For equal orientations,
\(\norm{Y_L+Y_R}^2=Q_{f,L}+Q_{f,R}+(\gamma_L+\gamma_R)/3\).
For mixed orientation it is
\(Q_{f,L}+Q_{f,R}+2\ip{b_L}{b_R}+2\gamma_{\rm near}/3\),
where \(\ip{b_L}{b_R}\leq5/12\).
Thus the equal/mixed bounds are \(17/6,11/3\) for two finite
owns, and \(3.077,3.91\) when a dotted own is present.
These bounds include unequal parameters and both mixed orientations.

If both ends are weak, pay \(\mu=.1\) on each foot.
The total individual anchor prices are bounded by
\(31/320,101/960,17/192\), all \(<.106<.134\).
For one weak end, pay \(\mu=.13\).
The near price is at most \(.13(31/32)<.126\), with
residual square \(<1.539\).
For far, the own-anchor price is at most \(.13(85/96)<.116\),
the separate opposite price is \(.13/10=.013\),
and its residual square is \(<1751/1500<1.168\).
Each fits its own available anchor; a bound on total price alone
would not suffice.

In all cases the payment is one completed square.
If the reference energy is \(E_0\), the residuals are \(Y_j\),
and each projection has been paid from nonnegative individual
strong-site coefficients, then
\[
 \ip y{w+\mu\sum_jY_j}\geq E_0+\mu\sum_j\Delta_j,\qquad
 \norm{w+\mu\sum_jY_j}^2=E_0+\mu^2\norm{\sum_jY_j}^2 .
\]
For two weak ends, adding the one \(I\) coordinate gives
\[
 e_A+\eta^2>
 .7986+.2(.358+.358)-.01(3.91)+.1164=1.0191.
\]
For one weak end the gain is \(>.067\):
\(.26(.358)-.0169(1.539)=.0670709\) near, and
\(.26(.336)-.0169(1.168)=.0676208\) far.
Thus \(e_A+\eta^2>.8284+.067+.1164=1.0118\).

If neither end is weak, use the known actual signed endpoint
coordinate \(>.51\) from \eqref{eq:B18A17interface}.
On its five-block the mean \(b_5\) satisfies \(.28<b_5<.28812\).
Adding \(\lambda(e_{\rm end}-P_5/5)\), with
\(\lambda=5(.51-b_5)/4\), spends only
\((.51-b_5)/(4\sqrt2)<.042<.134\) at that anchor.
It gains \((5/4)(.51-b_5)^2>.061538418\).
Consequently \(e_A+\eta^2>.858+.061538418+.1164>1\).
All weak-end counts are exhausted. The \(I\) energy was counted once
and every comparison used the same original farthest vertex.
This excludes the final orthogonal branch and hence every dotted
far-helper seventeenth cap.

\subsection{Near-three/own-six: no short opposite survives}
\label{app:B18near6}
The other left-contact terminal role has
\[
 g=\sqrt3/2,\quad
 c=Ap-\sqrt2gP_{17}-P_{16}/\sqrt2,\quad
 a=65/2+16\sqrt3,\quad Q=\sqrt2(17g+8),\quad Q^2=17a+9 .
\]
The inventory made old \(17\) strong, so there is the actual
isolated native-five head
\(c_H=A_Hp-(\phi/\sqrt2)e_1\), with
\(A_H^2=(\phi-1)/2\).
Use the common short table in Appendix~\ref{app:B18shorttable}.
It already excludes end-four, next-three and near end-three.
Far end-three has the same owns zero, three, four, six and parallel
nulls as before.

For any remaining finite far-three short own,
\(\phi/2\leq\delta<1\), its actual mutual magnitude with the head is
\[
 M_H(\delta)=A_H\sqrt{18\delta^2+\sqrt2\delta-\tfrac12}
             -\phi\delta/\sqrt2-\phi/2 .
\]
It increases strictly: the derivative of its first radical exceeds
\(\sqrt{18}\), and \(\sqrt{18}A_H>\phi/\sqrt2\).
At \(\delta=.8\), positive squares give time product \(>1.937\)
and spatial product \(<1.725\), so \(M_H>.212\).
At \(1\) the corresponding bounds are \(<2.419\) and
\(>1.9530878\), hence \(M_H<.4659122<1/2\).
Every finite remaining own has \(\delta>.8\), so
\(0<M_H<1/2\) contradicts the actual native gap.
The caps are distinct by their original contact sets.
Thus a far-three short opposite must have dotted own \(\delta>1\).

The following constants are deliberately coarse:
\[
\begin{gathered}
 60.212<a<60.2136,\quad 7.759<A<7.76,\quad32.13<Q<32.14,\\
 .1494<w:=9/a<.1495,\quad Q/A>4.141,\quad
 H:=A-G_{17}Q/17>.65,\quad
 0<K:=A\sqrt{18}-Q<.8 .
\end{gathered}
They follow by positive squares from \(1.732<\sqrt3<1.7321\);
for the last bound use
\((Q+.8)^2-18a=1.6Q+9.64-a>0\).
Positive actual clearing now uses the opposite inverse column:
\[
 \bar c=c+\sum_{i=1}^{16}(i/17)n_i
       =Ap-(Q/17)P_{17},\qquad \bar c^2=r=9/17 .
\]
This leaves the true short normal untouched.

We shall repeatedly use one quadratic seal.
For \(s>0,\rho>0,\rho^2>rs\), let
\(y_0=\sqrt{\rho^2/s-r}\).
If \(\rho<.8\delta\) and \(J>.31\delta\), then
\begin{equation}\label{eq:B18near6seal}
 H+(\rho/s)J-y_0
 >\frac{13}{20}+\frac{31}{80s}\rho^2-y_0
 =\frac{31}{80}(y_0-40/31)^2+\frac{8853}{42160}>0 .
\end{equation}

For an isolated short opposite \(d=Bp-\delta P_{18}\), put
\(\lambda=AB-Q\delta\geq0\).
Acuteness gives \(\delta^2\geq a/(a-9)>1\), so own is dotted.
The cleared Gram is
\(\left(\begin{smallmatrix}r&-\lambda\\-\lambda&1\end{smallmatrix}\right)\).
At \(\lambda^2=r\), \(\bar c+\lambda d\) is a positive
future all-wall feasible null.
For \(\lambda^2<r\), including \(\lambda=0\), the full positive
projection gives an actual finite point whose equations imply
\[
 k=w+\frac{Q\lambda}{a\delta}\geq w,\qquad k^2<w,\qquad
 S_{17}+kx_{18}\geq Q/A-\sqrt{w-k^2}.
\]
The two original sum bounds bound the left side by \(3.760+.07k\).
For \(k\geq w\), rationalization gives
\[
 \sqrt{w-w^2}-\sqrt{w-k^2}
 \geq\sqrt w\,(k-w)>.07(k-w)\quad(k>w).
\]
Hence the gap is minimized at \(k=w\). There it exceeds
\(4.141-.357-3.760-.07(.15)=.0135\), since
\(w-w^2<.1495-.1495^2<.357^2\).
Every subcritical mutual is impossible.
For \(\lambda^2>r\) the positive actual negative core is
\(\bar c+\lambda d\); its square is \(r-\lambda^2\),
selected products are \(r-\lambda^2,0\), its time is positive,
and all unknown signs hold termwise.
Its required pressure gap is
\(H+\lambda(B-G_{18}\delta)-\sqrt{\lambda^2-r}\).
Acuteness gives \(B/\delta\geq Q/A>4.141\), hence
\(J=B-G_{18}\delta>.31\delta\);
also \(\lambda<K\delta<.8\delta\).
Equation~\eqref{eq:B18near6seal} with \(s=1,\rho=\lambda\)
contradicts it on the entire supercritical range.

Finally consider dotted far-three short:
\(d=Bp-\delta P_{18}-e_1/\sqrt2\), \(\delta>1\).
Its mutual magnitude is
\(\lambda=AB-Q\delta-(g+1/2)\).
It is strictly positive on this whole range.
Its derivative is positive by the square comparison
\[
 18(a-9)\delta^2+\sqrt2(a-9)\delta+9a+\tfrac92>0 ,
\]
and at \(\delta=1\) the time/spatial square gap is
\(25/4+15g-\sqrt2(1+g)>16.50395\).
Clear \(d\) by the far inverse column to obtain
\[
 \bar d=Bp-(\delta+v_0)P_{17}-\delta e_{18},\quad
 \bar d^2=r,\quad
 \rho=-\ip{\bar c}{\bar d}=\lambda+1/34,\quad
 v_0=1/(17\sqrt2).
\]
The offset is now \(1/34\), not \(8/17\).
If \(\lambda=1/2\), then \(\rho=r\) and
\(\bar c+\bar d\) is a forbidden positive future feasible null.
The actual native gap then forces \(\lambda\geq1/\sqrt2\)
and \(\rho>r\).
Use \(\Xi=\bar c+(\rho/r)\bar d\).
It has square \(r-\rho^2/r<0\), positive time,
selected products that same negative square and zero, and zero
old \(A_{16}\) products. Every unselected original sign holds
termwise. The necessary same-\(p\) gap is
\[
 H+(\rho/r)J-\sqrt{\rho^2/r-r},\qquad
 J=B-G_{17}v_0-G_{18}\delta .
\]
The actual identity \(\rho=AB-Q(\delta+v_0)\) implies
\(J>.31\delta\). Moreover \(B<\sqrt{18}\delta+1/6\) and
\(Qv_0=g+8/17>A/6\), so \(\rho<K\delta<.8\delta\).
Equation~\eqref{eq:B18near6seal} with \(s=r\) excludes the
entire remaining range, including arbitrarily large dotted mutuals.
Thus no short opposite survives near-three/own-six.

\subsection{A central \(A_{16}\) lemma with a rank-two spectator}
\label{app:B18localA16}
The no-left terminal roles will lead to three types of actual
\(A_{16}+I_2(m)\) neighbors. We prove one local lemma with the
precise three comparison regimes:
\[
\begin{array}{c|c|c}
\text{spectator}&\text{comparison factor }\kappa&
 \text{conclusion on the same local \(A_{16}\) projection}\\ \hline
 I_2(3)&\kappa<67/50&e_A>1/2\\
 I_2(m),\ m\geq4&\kappa<3/2&e_A>7/16\\
 I_2(5)&\kappa<157/100&e_A>15/32 .
\end{array}
\]
Let the actual vertex be \(q\), write \(z=E_q(q-y)\), and suppose
every original chamber point has height at most \(\kappa E_q\).
Only the actual opposite \(b\) in central direction \(8\) is needed.
In zero-sum \(\mathbb R^{17}\),
\[
 \Omega_i=\sqrt2(P_i-iP_{17}/17),\quad
 D_i=2i(17-i)/17,\quad C_{ij}=2i(17-j)/17\ (i\leq j).
\]
Put \(D=D_8=144/17\), so \(2.91<\sqrt D<3\) and
\(\sqrt{D-1}>2.73\).
Write \(b=Hq-v_A-v_S\), own \(\gamma=-\ip b{m_8}\geq0\),
and whole-spectator cost \(Q=\norm{v_S}^2<1\).

The complete foreign alphabet follows from the retained dihedral
inverse. Two contacts make a nonpositive triangle.
For one finite contact \(\epsilon\) the cost is
\(\epsilon^2/\sin^2(\pi/m)\), and finiteness requires
\(\cos^2(\pi/m)+\epsilon^2<1\). Thus
\[
\begin{array}{c|c}
m& Q\\ \hline
3&0,\ 1/3,\ 2/3,\ \phi^2/3\\
4&0,\ 1/2\\
5&0,\ (\phi+2)/5\\
\geq6&0 .
\end{array}
\]
In particular the genuine \(H_3\) cost for \(m=5\) is
\((\phi+2)/5=(5+\sqrt5)/10>7/10\).
The whole block is cleared by its positive inverse on actual roots.

Deleting \(8\) retains \(A_7\) and \(A_8\).
The complete finite-white overlist is one three at
\(j=1,2,3,5,6,7,9,10,15,16\);
one four at \(j=1,7,9,16\);
or two end-threes, one from \(1,7\) and one from \(9,16\);
isolation is also allowed.
The left \(3,5\) positions are the \(E_8\) additions,
\(2,6,10,15\) are \(D\) additions. Two contacts on one path
make a cycle. A joined rank-sixteen tree is classical and has
no new internal four. The exceptional finite classification gives
no other extension of either retained full path.
Every nonempty white cost is at least \(7/16\).
Thus its retained Schur remainder is at most \(9/16\);
for \(m=5\) any white contact forces \(Q=0\).

Here is the complete own-zero boundary.
For an endpoint-three the time square is \(Q-9/17\), but
retention requires \(Q<9/16\). The displayed alphabets allow only
\(0,1/3,1/2\) under that bound, all too small.
For next-to-end three the retained remainder is at most \(1/4\),
hence \(Q=0\) and time square \(-2/17\).
For endpoint-four it is at most \(1/8\), hence \(Q=0\)
and time square \(-1/17\).
Every other internal-three has a positive affine principal:
\(E_8\) with arms \(1,2,5\) at global \(3,14\), or
\(E_7\) with arms \(1,3,3\) at \(4,\ldots,13\).
Their kernel coefficients are respectively center \(6\), arms
\(3;\ 4,2;\ 5,4,3,2,1\), and center \(4\), arms
\(2;\ 3,2,1;\ 3,2,1\).
Interior-four gives the \(4,3,3\) star with center \(\sqrt2\),
four-leaf \(1\), simple leaves \(1/\sqrt2\).
Two whites give an all-three cycle with coefficients one.
With no white contact the time square is \(Q-1<0\).
Every null is a positive actual combination, with positive cap
time, selected products zero and every other original product
nonpositive. Compactness excludes all of them.

For positive own, an own-three together with white-three closes an
actual all-three cycle through \(8\). For a pure far-three helper,
own-four gives the \(4,3,3\) star at \(8\), own-six gives the null
\(\gamma b+m_8+m_7/2\), and parallel gives \(b+m_8\).
The same own nulls will be used in isolated cases below.
The remaining ordinary edge ratios are obtained from the exact formula
\[
 R=\frac{C_8}{\sqrt D}
 =\frac{\gamma+\beta C_{8j}/D}
 {\sqrt{1-\beta^2(D_j-C_{8j}^2/D)-Q}}.
\]
For two whites add both numerator terms; their retained cross is zero.
We list bounds for every ordinary position, leaving only isolation
and pure far-three as separate cases:
\[
\begin{array}{c|c|c}
\text{white position}&\text{numerator lower bound}&
 \text{denominator upper bound before foreign cost}\\ \hline
\text{near three left/right}&.7+7/16,\ .7+4/9&3/4,\ \sqrt5/3\\
D\text{ left/right}&.7+1/8,\ .7+1/9&1/2,\ \sqrt2/3\\
E\text{ left}&.7+3/16&1/4\\
\text{four left/right}&.5+1/(8\sqrt2),\
 .5+1/(9\sqrt2)&1/\sqrt8,\ 1/3\\
\text{two threes}&.5+17/144&\sqrt{17}/12 .
\end{array}
\]
Smaller owns in these rows are the actual cycle/affine boundaries
already disposed of. For a far-three with \(Q>0\), the
\(m=3\) alphabet gives \(Q\geq1/3\); the residual is smaller
than \(.56^2\) and the numerator exceeds \(.864\).
For \(m\geq4\), positive \(Q\geq1/2\) makes this bound stronger.
Thus in the first two comparison regimes every ordinary row has
\(R>3/2\), and \(C_8>4.36\).
The directional floor is then
\[
 (C_8-\kappa)/\sqrt{C_8^2-1}>
 \begin{cases}(4.36-1.34)/4.244>.71,\\
 (4.36-1.5)/4.244>.67 .
 \end{cases}
\]
Its square exceeds \(1/2\) or \(7/16\), respectively.

For the third regime, white contact already forces \(Q=0\).
One near-three/own-four needs an actual mate, rather than an
incorrect null shortcut.
If \(v_A=\Omega_8/\sqrt2+\Omega_7/2\), delete actual \(7\);
the full finite Gram is \(A_6+B_{10}+I_2(5)\).
For \(\Omega_9/2\) instead, delete \(9\) and obtain
\(B_9+A_7+I_2(5)\).
Their positive actual projections are feasible against every original
wall. The exact edge data are
\[
\begin{array}{c|c|c}
j&\sigma_j&\Delta_j\\ \hline
7&1/2+9/(10\sqrt2)&1/10\\
9&1/2+8/(9\sqrt2)&1/9 .
\end{array}
\]
As \(\sqrt{D_j}>14/5,\ \sigma_j>9/8,\ \sqrt{\Delta_j}\leq1/3\),
their cosines exceed \(189/20>9\).
The new demand is \(>(9-1.57)/\sqrt{80}>.82\), more than sufficient.
An own-four with no neighboring single-three has the star null;
two-white or white-four roles use their ordinary bounds.
Every remaining white-three has own at least five, \(>.809\).
The same ratio formula gives:
near-three \(>(\phi/2+7/16)/(3/4)>1.662\);
\(D>(.809+1/9)/(.5)>1.84\);
\(E>(.809+3/16)/(.25)\);
far-four at least \(\sqrt2+1/4>1.664\) on the left and
\(3/2+1/(3\sqrt2)>1.7\) on the right;
two whites \(>(.7+17/144)/.344>2\).
Thus all these rows have \(R>1.66,\ C_8>4.83\), and demand
\[
 (C_8-1.57)/\sqrt{C_8^2-1}>
 (4.83-1.57)/4.726>11/16 .
\]
Its square is \(121/256>15/32\).

It remains to cover isolated \(A\) contacts, retaining every permitted
foreign cost. Put \(\tau=(1-Q)/\gamma^2\); then
\(H^2=\gamma^2(D-\tau)\).
Own three, four, six and parallel have the positive actual nulls
already listed, with own-three using the \(E_7\) star at \(8\).
For \(0<\tau<1\), positive clearing of the whole spectator gives
\[
 \Xi=b+w_S+\gamma\tau m_8,\quad
 U=\gamma(\Omega_8-\tau m_8),\quad
 \Xi^2=-\gamma^2\tau(1-\tau),\quad U^2<H^2 .
\]
It has zero cap/foreign products, negative own product, positive time
and every unknown sign termwise.
The comparison gives \(\ip{y_A}U\geq H-\kappa\gamma
\sqrt{\tau(1-\tau)}\).
If \(e_A\leq1/2\), Cauchy instead forces
\[
 (H-\ip{y_A}U)/\gamma
 >(1-1/\sqrt2)\sqrt{D-\tau}
 >.2928(2.73)=.799344>\kappa/2 ,
\]
a contradiction for every \(\kappa\leq1.57\).
Thus this whole continuous interval gives the stronger \(e_A>1/2\)
in all regimes. The boundary \(\tau=1\) is an actual null.

For \(\tau>1\) use the actual positive finite point
\[
 q_0=q+\frac{H}{\gamma^2(\tau-1)}(b+w_S+\gamma m_8).
\]
Its selected products vanish, its norm and \(q\) product are
\(-1-H^2/[\gamma^2(\tau-1)]\), and every unknown wall sign
holds by positive actual coefficients.
Normalization gives exactly
\[
 h_8\geq\sqrt{\frac{D-1}{D-\tau}}\,
          \{\sqrt{D-1}-\kappa\sqrt{\tau-1}\}.
\]
For \(Q>0\), the native gap gives
\(\tau-1\leq5-3\phi<(5/12)^2\).
For pure finite own at least eight use the sharper
\(\sqrt{\tau-1}\leq\tan(\pi/8)<.415\).
In the first regime the absolute loss in dual height satisfies
\[
 \sqrt D-h_8<.18+1.34(.415)=.7361<.74 .
\]
Since \(\sqrt D>2.91\), normalization gives
\[
 h_8/\sqrt D>1-.74/2.91>.74,
\]
and hence energy \(>.5476>1/2\);
the same bound covers the foreign case with its
\(\sqrt{5-3\phi}<.383<.415\).
In the second, \(h_8>2.73-1.5(.415)=2.1075\), whose normalized
value is \(>2/3\), giving energy \(>4/9>7/16\).
In the third \(h_8>2.73-1.57(.415)=2.07845>2.07\),
giving normalized value \(>.69\) and energy \(>.4761>15/32\).

Pure own-five has the actual \(H_4\) correction
\(\gamma(3m_8/2+m_7+m_6/2)\), with
\(\delta_5=(7-3\sqrt5)/16<1/36,\ H^2>4,\gamma>.8\).
Its actual finite point yields
\[
 h_8>\sqrt{D-\gamma^{-2}}-
       \kappa\sqrt{\delta_5}\sqrt{1+\delta_5/H^2}/\gamma .
\]
For \(\kappa\leq1.34\) its loss is
\(<25/76+1.34(1.004)/(6(.8))<.611<.74\),
so energy \(>.5476\).
For \(1.5\), the first radical is \(>2.5\) and the last square
root \(<1.01\), giving \(h_8>699/320\), enough for \(4/9\).
For \(1.57\), the first radical is \(>2.62\), the loss
\(<15857/48000\), and \(h_8>2.07\).
Pure seven has the actual core
\(b+(4\gamma/3)m_8+(2\gamma/3)m_7\), whose norm and cap
product are \(1-4\gamma^2/3<0\), with old products zero.
The cubic \(8u^3-4u^2-4u+1\) is increasing on \([.89,1]\);
its signs at \(.9,10/11\) give \(.9<\gamma<10/11\).
Thus \(\gamma^{-2}<100/81\) and
\(\sqrt{4/3-\gamma^{-2}}<9/25\). The pressure is
\[
 h_8>\sqrt{D-\gamma^{-2}}-\kappa\sqrt{4/3-\gamma^{-2}}.
\]
For the three regimes use respectively loss \(<.25+1.34(.36)<.74\),
first radical \(>2.6\), or first radical \(>2.68\).
The latter gives \(h_8>2.1148>2.07\).
All required isolated bounds follow.

Finally consider every pure far-end-three helper, including its
unbounded dotted range:
\[
 v_A=\gamma\Omega_8+\Omega_e/2,\quad e=1,16,\qquad
 H^2=D\gamma^2+C_{8e}\gamma-9/17>D\gamma^2,\quad
 C_{8e}\geq16/17,\quad \gamma>.8 .
\]
Own three, four, six and parallel have the actual cycle/star/chain
nulls above. For five, the actual correction
\(\gamma m_8+(2/3)m_e+(1/3)m_j\), where \(j=2,15\),
has residual \(\delta=2/3-\gamma^2\in(0,1/36)\).
The positive finite point gives
\(\ip{y_A}U>H-\kappa\sqrt\delta\sqrt{1+\delta/H^2}\),
with \(U^2=H^2+\delta\).
Even at \(\kappa=1.57\) the loss is \(<15857/60000<2/7\);
as \(H>2\), \(U^2<(10/9)H^2\), this gives
\(e_A>162/245>1/2\).
For finite owns at least seven let \(d_0=1-\gamma^2\in(0,.19)\).
The core \(b+\gamma m_8+2d_0m_e\) has square
\(-d_0(1-4d_0)\), zero cap/own products,
and helper product \(-1/2+2d_0<0\).
Its radical is at most \(1/4\), while \(H>2.5\);
the demand is \(>H-157/400>.84H\), giving \(e_A>441/625>1/2\).
For dotted \(\gamma>1\), the reciprocal core \(b+m_8/\gamma\)
has square \(-(1-\gamma^{-2})\).
Its radical divided by \(\gamma\) is at most \(1/2\).
As \(H>\gamma\sqrt D\), the demand/norm ratio is
\(>1-1.57/(2\sqrt D)>.73\), giving \(e_A>.5329>1/2\).
All cores have positive actual coefficients and time; all selected
products were checked, and all other original-wall products
are nonpositive termwise.
This finishes the three-regime local lemma without a new farthest
choice, a helper cutoff on arbitrary caps, or an omitted foreign cost.

\subsection{No-left caps: actual-neighbor recognition and all common cases}
\label{app:B18noleft}
The remaining seventeenth opposite has the form
\[
 c=Ap-\sqrt2gP_{17}-\delta P_{18},\qquad
 \delta=\cos(\pi/m),\quad 3\leq m\leq8,
\]
with the own table in Appendix~\ref{app:B18inventory}, and
\[
 A^2=34g^2+34\sqrt2g\delta+18\delta^2-1 .
\]
First distinguish actual common and distinct caps.
The full Gram on old \(A_{17}\) and \(c\) has Schur residual
\(1-17g^2/9\). If \(g=0,1/2,1/\sqrt2\), this is positive;
its positive full projection is the actual old short-edge endpoint.
The actual seventeenth opposite is therefore also the actual short
opposite, by uniqueness of the old edge endpoint.
For \(g=1/\sqrt2\) the reciprocal short cosine is
\(18\delta+17>9.81\), excluding these common-own-four roles.
For \(g\geq\phi/2\) the residual is negative, so the cap cannot
also close the short edge. This is the complete distinction for
the own values left by the table; no native product is applied to
a wall paired with itself.

At \(g=0,m=3\), the actual affine \(F_4\) path
\(n_{15},n_{16},n_{17},n_{18},c\) has positive kernel
\((1,2,3,2\sqrt2,\sqrt2)\), already excluded.
For \(g=\phi/2,m=4\), positive clearing of the whole old
\(A_{17}\) gives the actual negative core
\[
 \Xi=Ap-sP_{18},\quad
 A^2=(51\phi+33)/2,\quad s=\sqrt2(18+17\phi)/36,\quad
 \Xi^2=-(17\phi-19)/36 .
\]
Its selected old products vanish, its cap product is its negative
square, its time is positive, and every other original sign holds
termwise. Its pressure implies
\[
 S_{18}>\{A-\sqrt{(17\phi-19)/36}\}/s>3.97,
\]
using \(A>7.599\), radical \(<.487\), \(s<1.79\).
This contradicts the unconditional sum bound.
Thus the only distinct-cap family has \(m=3\) and
\(\phi/2\leq g<1.264\).

For each of these remaining families let \(q\) be the actual
seventeenth endpoint. Its complete actual forest is
\(A_{16}(n_1,\ldots,n_{16})+I_2(m)(c,n_{18})\).
The cap and short root have their original finite mutual
\(-\cos(\pi/m)\); the full Gram is positive.
Its full positive actual projection checks every unknown wall.
Put
\[
 C=\cosh d(p,q),\quad s=\sqrt{C^2-1},\quad
 t=S_{17}/\sqrt{17},\quad
 q=Cp-sP_{17}/\sqrt{17},\quad E_q/E=C-st .
\]
The reverse unit tangent belongs wholly to the dihedral component.
Its pressure is
\begin{equation}\label{eq:B18reverse}
 \frac{\beta-t}{1-\beta t},\qquad \beta=s/C .
\end{equation}
It increases in \(\beta\) and decreases in \(t\).
Thus the reverse-direction energy and the \(A_{16}\) energy
from Appendix~\ref{app:B18localA16} are orthogonal projections
of the same local vector.

For the distinct family \(m=3,g\geq\phi/2\),
\[
 C=\sqrt{136/3}\{g+1/(2\sqrt2)\}>7.82,\qquad t<.912.
\]
The function \(C-.912\sqrt{C^2-1}\) increases here, so
\[
 E_q/E>7.82-.912(7.756)=.746528>50/67 .
\]
Also \(\beta>.9917\), and \eqref{eq:B18reverse} is \(>.83\).
The local \(A_{16}+I_2(3)\) lemma gives \(e_A>1/2\), while
the actual dihedral projection has energy \(>.6889\).
Their sum exceeds one, excluding every distinct cap.

Of the seven common cases left by the own table, five are
\(g=0,m=6,7,8\) or \(g=1/2,m=4,5\).
Their exact forward cosine is
\[
 C=\sqrt{34}\,\frac{g+\delta/\sqrt2}{\sqrt{1-\delta^2}} .
\]
For \(g=0,m\geq6\), \(C\geq\sqrt{51}>7.1\);
for \(g=1/2,m\geq4\), \(C\geq\sqrt{68}\).
Again \(t<.912\), hence
\[
 E_q/E>7.1-.912(7.03)=.68864>2/3 .
\]
Now \(\beta>.99\), so the reverse pressure is
\(>(.99-.912)/(1-.99(.912))>.8\).
The second local regime gives \(e_A>7/16\); the dihedral
projection costs \(>.64\). Their sum is \(>431/400>1\).
This excludes all five common cases.

For the sixth case, \(g=0,m=5\), we first improve the original
comparison without adding a spent head energy.
The actual cap is \(c=Ap-(\phi/2)P_{18}\), with
\(A^2=(9\phi+7)/2\).
If old \(17\) were strong the actual head would be native five.
The distinct head and cap would then have mutual magnitude
\[
 M=\sqrt{7\phi+2}/2-(\phi+1)/(2\sqrt2).
\]
Rational squares from \(1.618<\phi<1.6181\) and
\(1.4142<\sqrt2<1.4143\) give \(.899<M<.900\).
But this lies strictly between the native-six and native-seven
values, a forbidden actual mutual.
Therefore \(S_{17}\leq f_{17}\), and \(t<.901\).
For its actual \(A_{16}+I_2(5)\) neighbor,
\[
 C^2=\frac{17(\phi+1)}{3-\phi}>32.2>5.67^2,\qquad
 \beta>.984,\qquad
 E_q/E>5.67-.901(5.582)=.640618>100/157 .
\]
The reverse pressure is
\(>(.984-.901)/(1-.984(.901))>.73\).
The third local regime gives \(e_A>15/32\), while the
orthogonal dihedral energy is \(>.5329\).
The sum exceeds \(20033/20000>1\).
The head gap was used solely to establish old weakness; no
head payment was added to this neighboring energy.

The final common case is \(g=1/2,m=3\).
Here \(c=Ap-P_{17}/\sqrt2-P_{18}/2\), with
\(A^2=12+(17/2)\sqrt2\).
The full positive old-\(A_{17}\)-plus-\(c\) projection is the
actual short endpoint \(q\), and its whole incident diagram is
\(A_{18}\), not the \(A_{16}+A_2\) diagram of the other endpoint.
Its Schur residual is \(19/36\), and
\[
 C=\frac{17+9\sqrt2}{\sqrt{19}},\quad
 C^2=(451+306\sqrt2)/19>6.82^2,\quad
 t=S_{18}/\sqrt{18}<.903 .
\]
Consequently
\[
 E_q/E>6.82-.903(6.747)=.727459>8/11 .
\]
The reverse unit tangent is the positive normalized dual for the
new \(A_{18}\) endpoint \(c\). With \(\beta>.989\),
its pressure is
\[
 \ip{y}{u_{18}}
 >\frac{.989-.903}{1-.989(.903)}>.8>.78 .
\]
Every original chamber point has height at most \(E<(11/8)E_q\).
The exact hypotheses of Lemma~\ref{lem:localA18} therefore hold
at this actual whole-\(A_{18}\) vertex.
That lemma is local and does not require \(q\) farthest.
It excludes this last common case. Its endpoint threshold concerns
the normalized dual direction; it is not a threshold on a raw
signed coordinate. All original wall signs here come from the
full positive finite projection.

\subsection{The terminal composition}
\label{app:B18conclusion}
The exhaustive seventeenth inventory has three disjoint branches.
A far-left-three helper has finite own five or at least seven,
or dotted own \(1<g<4/3\).
Appendix~\ref{app:B18finitefar} excludes every finite own, and
Appendices~\ref{app:B18dottedfar}--\ref{app:B18localA17}
exclude every dotted own, every actual short-white profile
and every zero/native/parallel/unbounded dotted mutual.
The near-left-three branch has exactly own six and is excluded
by Appendix~\ref{app:B18near6}.
The no-left/short-contact branch is exhausted by the distinct/common
recognition and the seven common cases above.
Every other white role and every isolated seventeenth opposite
was removed before this split.
It follows that an original farthest whole-\(B_{18}\) vertex
cannot occur.

All arguments are geometric. Literal finite/affine classification,
compact acute simplicity, inverse positivity and the actual unit-wall
angle alphabet are the external mathematical inputs.
The finite tables here are explicit rational/algebraic inequalities,
not new enumeration outputs or runtime receipts.
The one original farthest vector is retained throughout each original
reference argument. A neighboring vector is introduced only with its
derived comparison factor, and all its orthogonal energies are taken
at that same actual neighbor.
No geometric Lean formalization, new script PASS or complete
compact-dimension theorem is asserted by this unintegrated fragment.
The remaining whole-\(D_{18}\), higher-rank and forest consumers
belong to the main manuscript's separate dependency closure.





